\documentclass{article}
\usepackage{graphicx}
\usepackage{amsmath}

\title{VermiHeatLoss}
\author{Ava Tob\'on}
\date{March 2026}

\begin{document}

\subsection{Overall Heat-Transfer Paths and Coupled Energy Balances}

The MATLAB model represents the heater and vermicomposter as a set of coupled
heat-transfer paths rather than a single series resistance chain. The
principal paths are:

\begin{enumerate}
\item Electrical heating in the nichrome wire, with convection from the wire
      to the heater-tube airflow.
\item Loss from the heater-tube airflow to ambient through the heater tube
      wall and insulation.
\item Useful heat transfer from the aeration airflow to the substrate by
      \begin{itemize}
      \item transfer through the aeration tube wall,
      \item direct warm-air release into the bed through the perforations, and
      \item subtraction of a latent heat sink associated with evaporation from
            the wetted substrate.
      \end{itemize}
\item Heat loss from the bed to ambient through the vessel wall.
\end{enumerate}

\subsubsection{Wire-to-Air Convection}

The wire-to-air convective coefficient is computed as
\begin{equation}
h_{\text{wire}}=\frac{Nu_{\text{wire}}\,k}{D_{\text{wire}}}
\end{equation}
where $Nu_{\text{wire}}$ is obtained from the Churchill--Bernstein
cylinder-in-crossflow correlation [R1]:
\begin{equation}
Nu_{\text{wire}}
=
0.3+
\frac{0.62\,Re_{\text{wire}}^{1/2}Pr^{1/3}}
{\left[1+\left(0.4/Pr\right)^{2/3}\right]^{1/4}}
\left[1+\left(\frac{Re_{\text{wire}}}{282000}\right)^{5/8}\right]^{4/5}.
\end{equation}
Equivalently, matching the code implementation,
\begin{equation}
Nu_{\text{wire}}
=
0.3+
\left(
\frac{0.62\,Re_{\text{wire}}^{1/2}Pr^{1/3}}
{\left[1+\left(0.4/Pr\right)^{2/3}\right]^{1/4}}
\right)
\left[1+\left(\frac{Re_{\text{wire}}}{282000}\right)^{5/8}\right]^{4/5}.
\end{equation}
The same Churchill--Bernstein form is also used for the perforation
hole-scale crossflow convection in the moisture model.

This coefficient is used to estimate the local wire-to-air temperature rise:
\begin{equation}
\Delta T_{\text{wire}}=\frac{q_{\text{gen}}}{h_{\text{wire}}A_{\text{wire,local}}}
\end{equation}
so that
\begin{equation}
T_{\text{wire}}=T_{\text{air}}+\Delta T_{\text{wire}}.
\end{equation}

This term affects wire temperature prediction, but it is not part of an
overall $U$ value.

\subsubsection{Heater-Tube Loss to Ambient}

The internal convection coefficient for the heater tube is
\begin{equation}
h_{\text{inside,heater}}=\frac{Nu_{\text{tube}}\,k}{D_i}.
\end{equation}
The implemented internal-tube Nusselt correlation is piecewise. For laminar
developing flow the code uses a Hausen-type entrance-length correction [R5]:
\begin{equation}
Gz = Re\,Pr\,\frac{D_i}{L},
\end{equation}
\begin{equation}
Nu_{\text{tube}}
=
3.66+\frac{0.0668\,Gz}{1+0.04\,Gz^{2/3}},
\qquad Re<2300.
\end{equation}
For turbulent flow the code uses Gnielinski [R2]:
\begin{equation}
Nu_{\text{tube}}
=
\frac{(f/8)(Re-1000)Pr}
{1+12.7(f/8)^{1/2}(Pr^{2/3}-1)},
\qquad Re\ge 2300,
\end{equation}
with the lower bound
\begin{equation}
Nu_{\text{tube}} \ge 3.66.
\end{equation}
Here $f$ is the Churchill friction factor [R3] used consistently in both the
pressure-drop and turbulent heat-transfer calculations.

The overall heat-transfer coefficient from the heater-tube air to ambient is
modeled as a series resistance network:
\begin{equation}
U_{\text{heater}}
=
\left(
\frac{1}{h_{\text{inside,heater}}}
+\frac{OD-ID}{2k_{\text{wall}}}
+\frac{t_{\text{ins}}}{k_{\text{ins}}}
+\frac{1}{h_{\text{outside,heater}}}
\right)^{-1}.
\end{equation}

The corresponding heat loss from the heater section is
\begin{equation}
q_{\text{loss,heater}}
=
U_{\text{heater}}\,P_{\text{out}}\,\Delta x\,
\left(T_{\text{air}}-T_{\text{env}}\right).
\end{equation}

This loss is coupled to the heater air-energy balance through
\begin{equation}
q_{\text{to air}}=q_{\text{gen}}-q_{\text{loss,heater}}
\end{equation}
and
\begin{equation}
\dot m c_p\,dT=q_{\text{to air}}.
\end{equation}

Thus, $U_{\text{heater}}$ determines how much electrical power remains in the
airflow versus being lost directly to ambient.

\subsubsection{Aeration-Tube Heat Transfer to the Bed}

The internal aeration-tube convection coefficient is
\begin{equation}
h_{\text{inside,aer}}=\frac{Nu_{\text{tube}}\,k}{D_i}
\end{equation}
and the external bed-side coefficient is taken as
\begin{equation}
h_{\text{outside,aer}}=h_{\text{bed}}.
\end{equation}
The same internal-tube $Nu_{\text{tube}}$ branch logic given above is used
for the aeration tube as for the heater tube: Hausen-type laminar developing
flow for $Re<2300$ and Gnielinski for $Re\ge 2300$.

The overall coefficient from aeration air to bed through the aeration tube
wall is
\begin{equation}
U_{\text{aer}}
=
\left(
\frac{1}{h_{\text{inside,aer}}}
+\frac{OD-ID}{2k_{\text{wall}}}
+\frac{1}{h_{\text{outside,aer}}}
\right)^{-1}.
\end{equation}

Wall-mediated heat transfer from air to bed is then
\begin{equation}
q_{\text{wall}}
=
U_{\text{aer}}\,P_{\text{out}}\,\Delta x\,
\left(T_{\text{air}}-T_{\text{bed}}\right).
\end{equation}

The model also includes direct jet heating from released warm air:
\begin{equation}
q_{\text{jet}}
=
\dot m_{\text{release}}\,c_p\,
\left(T_{\text{air}}-T_{\text{bed}}\right).
\end{equation}

In the current implementation, evaporation from the wetted substrate is
modeled explicitly as a latent sink:
\begin{equation}
q_{\text{evap}}
=
\dot m_{\text{evap}}\,h_{fg}.
\end{equation}

Hence the branch heat delivered to the bed is
\begin{equation}
Q_{\text{bed,branch}}=\sum \left(q_{\text{wall}}+q_{\text{jet}}-q_{\text{evap}}\right)
\end{equation}
and the total heat delivered by all parallel tubes is
\begin{equation}
Q_{\text{bed,total}}=N_{\text{tubes}}\,Q_{\text{bed,branch}}.
\end{equation}

Therefore, $U_{\text{aer}}$ sets the wall-transfer contribution to useful bed
heating, $q_{\text{jet}}$ adds a parallel direct-heating path, and
$q_{\text{evap}}$ reduces the net useful heat delivered to the substrate.

The bed-side coefficient $h_{\text{bed}}$ remains a \emph{sensible}
wall-transfer coefficient. Latent evaporation is modeled separately through
$q_{\text{evap}}$ and is not absorbed into $h_{\text{bed}}$.

\subsubsection{Vessel-Wall Heat Loss to Ambient}

For the bin exterior, the natural-convection coefficient is obtained from
\begin{equation}
Ra=\frac{g\beta\Delta T L^3}{\nu\alpha},
\end{equation}
followed by the Churchill--Chu vertical-plate relation [R4],
\begin{equation}
Nu
=
\left[
0.825+
\frac{0.387\,Ra^{1/6}}
{\left(1+(0.492/Pr)^{9/16}\right)^{8/27}}
\right]^2,
\end{equation}
and then
\begin{equation}
h_{\text{ext}}=\frac{Nu\,k}{L}.
\end{equation}

The overall vessel-wall coefficient is
\begin{equation}
U_{\text{wall}}
=
\left(
\frac{1}{h_{\text{inside,bin}}}
+\frac{t_{\text{sheet}}}{k_{\text{sheet}}}
+\frac{t_{\text{ins}}}{k_{\text{ins}}}
+\frac{1}{h_{\text{ext}}}
\right)^{-1}.
\end{equation}

The total bed-to-ambient heat loss is
\begin{equation}
Q_{\text{loss}}
=
U_{\text{wall}}A_{\text{total}}
\left(T_{\text{bin}}-T_{\infty}\right).
\end{equation}

The same $U_{\text{wall}}$ is used to construct the nodewise ambient-loss
terms:
\begin{equation}
UA_{\text{bottom}}=U_{\text{wall}}A_{\text{bottom,node}},
\qquad
UA_{\text{top}}=U_{\text{wall}}A_{\text{top,node}}.
\end{equation}

Thus, $U_{\text{wall}}$ governs the final loss of stored bed heat to ambient.

\subsubsection{Coupling Through the Bed Energy Balances}

The bed is represented as two lumped thermal nodes: a bottom node and a top
node. Heat is added to these nodes from the aeration system, exchanged
internally between the two bed regions, and lost to ambient through the
vessel boundaries.

\paragraph{Heat input split between bottom and top nodes}

The total useful heat delivered from the aeration system is
\begin{equation}
Q_{\text{bed,total}} = N_{\text{tubes}} Q_{\text{bed,branch}}.
\end{equation}

This heat is then partitioned between the bottom and top bed nodes using the
parameter $f_b$:
\begin{equation}
Q_b = f_b\,Q_{\text{bed,total}},
\qquad
Q_t = (1-f_b)\,Q_{\text{bed,total}}.
\end{equation}

Here, $f_b$ is an assumed heat-input split parameter that represents the
heater tube proximity to the bottom of the bed. It does \emph{not} represent
ambient heat loss.

\paragraph{Internal coupling between bottom and top bed regions}

The two nodes are thermally coupled through an internal conduction-like term:
\begin{equation}
UA_{\text{internal}}
=
k_{\text{bed}}
\frac{A_{\text{internal}}}{\delta_{\text{internal}}},
\end{equation}
where
\begin{equation}
A_{\text{internal}} = LW
\end{equation}
is the horizontal interface area between the lower and upper bed regions, and
\begin{equation}
\delta_{\text{internal}} \approx \frac{H}{2}
\end{equation}
is the assumed characteristic distance between the node centroids.

This term allows heat to move between the bottom and top nodes:
\begin{equation}
Q_{\text{internal}} = UA_{\text{internal}}(T_t-T_b).
\end{equation}

\paragraph{Ambient heat loss areas for each node}

The bottom node loses heat through the bottom surface and half of the
side-wall area:
\begin{equation}
A_{\text{bottom,node}} = A_{\text{bottom}} + \frac{1}{2}A_{\text{side}},
\end{equation}
where
\begin{equation}
A_{\text{bottom}} = LW
\end{equation}
and
\begin{equation}
A_{\text{side}} = 2HL + 2HW.
\end{equation}

Thus the bottom-node ambient loss coefficient is
\begin{equation}
UA_{\text{bottom}} = U_{\text{wall}} A_{\text{bottom,node}}.
\end{equation}

For the top node, the treatment depends on whether the bin top is covered or
open.

If the top is treated as covered,
\begin{equation}
UA_{\text{top}}
=
U_{\text{wall}}
\left(
A_{\text{top}}+\frac{1}{2}A_{\text{side}}
\right),
\end{equation}
where
\begin{equation}
A_{\text{top}} = LW.
\end{equation}

If the top is treated as open,
\begin{equation}
UA_{\text{top}}
=
U_{\text{wall}}
\left(
\frac{1}{2}A_{\text{side}}
\right)
+
h_{\text{top}}A_{\text{top}}.
\end{equation}

Therefore, with an open top:
\begin{itemize}
\item side-wall losses are still present and are split equally between the
      two nodes, and
\item the top surface itself loses heat directly to ambient by convection
      rather than through the insulated wall stack.
\end{itemize}

\paragraph{Steady two-node balance}

Under steady conditions, the node temperatures $T_b$ and $T_t$ satisfy
\begin{equation}
(UA_{\text{bottom}}+UA_{\text{internal}})T_b
-
UA_{\text{internal}}T_t
=
Q_b + UA_{\text{bottom}}T_{\text{amb}}
\end{equation}
and
\begin{equation}
-UA_{\text{internal}}T_b
+
(UA_{\text{top}}+UA_{\text{internal}})T_t
=
Q_t + UA_{\text{top}}T_{\text{amb}}.
\end{equation}

These equations show that each node is heated by its assigned fraction of the
aeration heat input, cooled by ambient loss through its assigned external
area, and coupled to the other node through internal heat exchange.

\paragraph{Transient two-node balance}

In the transient model, the node temperatures evolve according to
\begin{equation}
\frac{dT_b}{dt}
=
\frac{
Q_b
+
UA_{\text{internal}}(T_t-T_b)
-
UA_{\text{bottom}}(T_b-T_{\text{amb}})
}{
C_b
}
\end{equation}
and
\begin{equation}
\frac{dT_t}{dt}
=
\frac{
Q_t
+
UA_{\text{internal}}(T_b-T_t)
-
UA_{\text{top}}(T_t-T_{\text{amb}})
}{
C_t
}.
\end{equation}

Here, $C_b$ and $C_t$ are the thermal capacitances of the bottom and top
nodes, respectively.

\paragraph{Interpretation}

The two-node bed model therefore contains three distinct mechanisms:

\begin{enumerate}
\item \textbf{Heater input distribution:} $f_b$ determines how the delivered
      aeration heat is partitioned between the bottom and top nodes.
\item \textbf{Internal redistribution:} $UA_{\text{internal}}$ allows heat to
      move between the two bed regions.
\item \textbf{Ambient loss:} $UA_{\text{bottom}}$ and $UA_{\text{top}}$
      determine how each node loses heat to the surroundings through the
      bottom, side walls, and top surface.
\end{enumerate}

This means that the top node may cool faster than the bottom node even if it
receives less direct heater input, particularly when the top is modeled as
open to ambient.

\subsubsection{Summary of Coupling}

The overall coupling structure is therefore:

\begin{equation}
\text{wire power}
\rightarrow
\text{heater air}
\rightarrow
\text{bed heat input}
\rightarrow
\text{bed storage}
\rightarrow
\text{ambient loss}.
\end{equation}

More specifically:
\begin{itemize}
\item $h_{\text{wire}}$ determines wire temperature relative to the heater
      air.
\item $U_{\text{heater}}$ determines how much heater-air energy is lost
      directly to ambient.
\item $U_{\text{aer}}$ determines how much aeration-air energy crosses the
      aeration tube wall into the bed.
\item $Q_{\text{bed,total}}$ combines wall transfer, jet heating, and
      evaporative cooling into the bed-energy input.
\item $U_{\text{wall}}$ determines how quickly the bed loses heat through the
      vessel to ambient.
\end{itemize}

\subsubsection{Meaning and Calculation of $T_{\text{air}}$ and $P_{\text{out}}$}

In the model, $T_{\text{air}}$ denotes the temperature of the flowing air
stream inside the heater tube or aeration tube, depending on the section
being analyzed. It is not a fixed temperature; rather, it is updated
segment-by-segment from an energy balance.

\paragraph{Heater-tube air temperature}

At the inlet to the heater tube, the air is assumed to enter at ambient
greenhouse temperature:
\begin{equation}
T_{\text{air,in}} = T_{\text{env}}.
\end{equation}

As the air moves through the heater, it gains energy from the electrically
heated wire and loses some energy to the surroundings through the heater tube
wall and insulation. For each axial segment, the net heat transferred to the
air is
\begin{equation}
q_{\text{to air}} = q_{\text{gen}} - q_{\text{loss,heater}},
\end{equation}
where
\begin{equation}
q_{\text{gen}} = \frac{P_{\text{branch}}}{L}\,\Delta x
\end{equation}
is the local rate of electrical heat generation assigned to the segment, and
\begin{equation}
q_{\text{loss,heater}} =
U_{\text{heater}}\,P_{\text{out}}\,\Delta x\,
\left(T_{\text{air}}-T_{\text{env}}\right)
\end{equation}
is the local heat loss from the heater air to ambient.

The air temperature is then advanced according to
\begin{equation}
T_{\text{air,new}}
=
T_{\text{air,old}}
+
\frac{q_{\text{to air}}}{\dot m c_p}.
\end{equation}

Thus, in the heater tube, $T_{\text{air}}$ is the local heated-air
temperature resulting from the balance between wire heating and heat loss to
the surroundings.

\paragraph{Aeration-tube air temperature}

The air entering the aeration tube is the outlet air from the heater tube:
\begin{equation}
T_{\text{air,in,aer}} = T_{\text{air,out,heater}}.
\end{equation}

As the air flows through the aeration tube, it transfers heat to the
substrate by three coupled mechanisms:
\begin{enumerate}
\item wall-mediated transfer through the aeration tube wall,
\item direct sensible heat release with the discharged warm air, and
\item latent heat consumption by evaporation from the wetted substrate.
\end{enumerate}

The wall-transfer term is
\begin{equation}
q_{\text{wall}}
=
U_{\text{aer}}\,P_{\text{out}}\,\Delta x\,
\left(T_{\text{air}}-T_{\text{bed}}\right),
\end{equation}
the direct jet term is
\begin{equation}
q_{\text{jet}}
=
\dot m_{\text{release}}\,c_p\,
\left(T_{\text{air}}-T_{\text{bed}}\right),
\end{equation}
and the evaporation term is
\begin{equation}
q_{\text{evap}}=\dot m_{\text{evap}}\,h_{fg}.
\end{equation}

In the present implementation, the air temperature along the aeration tube is
updated using the wall-transfer term:
\begin{equation}
T_{\text{air,new}}
=
T_{\text{air,old}}
-
\frac{q_{\text{wall}}}{\dot m_{\text{remaining}} c_p}.
\end{equation}

Therefore, in the aeration section, $T_{\text{air}}$ is the local air
temperature after heater exit, reduced progressively as sensible heat is
transferred from the air stream to the bed. The net bed heat delivered at the
same segment is then corrected by subtracting the latent evaporation sink.

\paragraph{Meaning of $P_{\text{out}}$}

The quantity $P_{\text{out}}$ is the \emph{outside perimeter} of the
cylindrical tube used to convert an axial segment length into an external
surface area. For a tube of outside diameter $D_{\text{out}}$,
\begin{equation}
P_{\text{out}} = \pi D_{\text{out}}.
\end{equation}

Hence,
\begin{equation}
P_{\text{out}}\,\Delta x
\end{equation}
is the outside surface area of one tube segment of length $\Delta x$.

This is why the wall heat-transfer terms are written in the form
\begin{equation}
q = U\,P_{\text{out}}\,\Delta x\,\Delta T:
\end{equation}
the product $P_{\text{out}}\Delta x$ supplies the heat-transfer area, and
$U\Delta T$ supplies the heat flux per unit area.

\paragraph{Summary}

In summary:
\begin{itemize}
\item $T_{\text{air}}$ is the local temperature of the moving air stream.
\item In the heater tube, it increases due to electrical heating minus
      ambient losses.
\item In the aeration tube, it decreases as sensible heat is transferred to
      the substrate.
\item The net useful bed heating in the aeration tube is further reduced by
      the latent evaporation term $q_{\text{evap}}$.
\item $P_{\text{out}}$ is the tube outside perimeter, so that
      $P_{\text{out}}\Delta x$ is the external area of each tube segment.
\end{itemize}

\paragraph{$h_{\text{top}}$ for an uncovered top surface}

When the bin top is modeled as open to ambient, the top-surface heat loss is
represented by a natural-convection coefficient $h_{\text{top}}$. This
coefficient describes free convection between the exposed upper bed surface
and the surrounding air.

It is defined by
\begin{equation}
h_{\text{top}}=\frac{Nu\,k}{L_c},
\end{equation}
where the characteristic length is taken as
\begin{equation}
L_c=\frac{A_{\text{top}}}{P_{\text{top}}},
\end{equation}
with
\begin{equation}
A_{\text{top}} = LW,
\qquad
P_{\text{top}} = 2(L+W).
\end{equation}

The Rayleigh number for the open top is
\begin{equation}
Ra=\frac{g\beta\Delta T L_c^3}{\nu\alpha}.
\end{equation}

Using the standard McAdams-style correlations for free convection from
horizontal plates:

For a \emph{heated upward-facing surface} (or equivalently a cooled
downward-facing surface),
\begin{equation}
Nu = 0.54\,Ra^{1/4},
\qquad
Ra \le 10^{7},
\end{equation}
and
\begin{equation}
Nu = 0.15\,Ra^{1/3},
\qquad
Ra > 10^{7}.
\end{equation}

For a \emph{heated downward-facing surface} (or equivalently a cooled
upward-facing surface),
\begin{equation}
Nu = 0.27\,Ra^{1/4},
\end{equation}

Since the uncovered compost surface is warmer than the ambient air during
heating operation, the physically relevant case is the heated upward-facing
surface. Therefore, the top-node ambient heat loss is modeled as
\begin{equation}
Q_{\text{top,ambient}}
=
h_{\text{top}}A_{\text{top}}(T_t-T_{\text{amb}}).
\end{equation}

In the two-node bed model, this contributes to the top-node ambient-loss
coefficient as
\begin{equation}
UA_{\text{top}}
=
U_{\text{wall}}\left(\frac{1}{2}A_{\text{side}}\right)
+
h_{\text{top}}A_{\text{top}},
\qquad \text{for an open top.}
\end{equation}

Thus, when the top is uncovered, the top-node heat loss consists of:
\begin{itemize}
\item side-wall loss through the vessel wall, represented by
\[
U_{\text{wall}}\left(\frac{1}{2}A_{\text{side}}\right),
\]
and
\item direct natural-convection loss from the exposed top surface,
      represented by
\[
h_{\text{top}}A_{\text{top}}.
\]
\end{itemize}

If the top is instead treated as covered, the model does not use
$h_{\text{top}}$; it applies the covered-wall coefficient $U_{\text{wall}}$
over the top area as well.

\paragraph{Total vessel heat loss}

The total vessel heat loss is the sum of heat losses from the bottom, sides,
and top surfaces to the ambient air.

\subparagraph{Closed-top vessel}

If the top is treated as covered, then the same overall wall coefficient
$U_{\text{wall}}$ applies to all external surfaces. The total vessel heat
loss is

\begin{equation}
Q_{\text{loss,closed}}
=
U_{\text{wall}}A_{\text{total}}(T_{\text{bin}}-T_{\infty}),
\end{equation}
where
\begin{equation}
A_{\text{total}} = A_{\text{side}} + A_{\text{bottom}} + A_{\text{top}}.
\end{equation}

Since
\begin{equation}
A_{\text{side}} = 2HL + 2HW,
\qquad
A_{\text{bottom}} = LW,
\qquad
A_{\text{top}} = LW,
\end{equation}
this may also be written as
\begin{equation}
Q_{\text{loss,closed}}
=
U_{\text{wall}}
\left(
2HL + 2HW + 2LW
\right)
(T_{\text{bin}}-T_{\infty}).
\end{equation}

\subparagraph{Open-top vessel}

If the top is treated as open, then the sides and bottom still lose heat
through the insulated vessel wall with coefficient $U_{\text{wall}}$, while
the exposed top loses heat directly to ambient by convection with coefficient
$h_{\text{top}}$. The total vessel heat loss is therefore

\begin{equation}
Q_{\text{loss,open}}
=
U_{\text{wall}}(A_{\text{side}}+A_{\text{bottom}})(T_{\text{bin}}-T_{\infty})
+
h_{\text{top}}A_{\text{top}}(T_{\text{bin}}-T_{\infty}).
\end{equation}

Substituting the areas gives
\begin{equation}
Q_{\text{loss,open}}
=
\left[
U_{\text{wall}}(2HL+2HW+LW)
+
h_{\text{top}}(LW)
\right]
(T_{\text{bin}}-T_{\infty}).
\end{equation}

\subparagraph{Equivalent overall conductance}

It is often convenient to write these in terms of an overall conductance $UA$.

For a closed top,
\begin{equation}
(UA)_{\text{closed}}
=
U_{\text{wall}}A_{\text{total}}.
\end{equation}

For an open top,
\begin{equation}
(UA)_{\text{open}}
=
U_{\text{wall}}(A_{\text{side}}+A_{\text{bottom}})
+
h_{\text{top}}A_{\text{top}}.
\end{equation}

Hence,
\begin{equation}
Q_{\text{loss}} = (UA)(T_{\text{bin}}-T_{\infty}),
\end{equation}
with the appropriate form of $UA$ selected according to whether the vessel
top is modeled as closed or open.

\subsection{Covered Top with a Small Vent Hole}

In addition to the limiting cases of a fully covered top and a fully open
top, the thermal model includes an intermediate case in which the bin top is
covered except for a small circular vent opening that allows the injected air
to escape. Let the total top area be
\begin{equation}
    A_{\mathrm{top}} = L W,
\end{equation}
and let the vent consist of $N_{\mathrm{vent}}$ identical circular holes of
diameter $d_{\mathrm{vent}}$. The total open vent area is then
\begin{equation}
    A_{\mathrm{vent}} = N_{\mathrm{vent}} \frac{\pi d_{\mathrm{vent}}^2}{4},
\end{equation}
with the covered portion of the top given by
\begin{equation}
    A_{\mathrm{top,cov}} = A_{\mathrm{top}} - A_{\mathrm{vent}}.
\end{equation}

The covered portion of the top is treated in the same manner as the insulated
vessel wall, using the overall wall coefficient $U_{\mathrm{wall}}$. The vent
opening is treated as a locally open horizontal surface that exchanges heat
directly with the greenhouse air by natural convection. For the vent opening,
the characteristic length is taken as
\begin{equation}
    L_{c,\mathrm{vent}} = \frac{A_{\mathrm{vent}}}{P_{\mathrm{vent}}},
\end{equation}
where the total vent perimeter is
\begin{equation}
    P_{\mathrm{vent}} = N_{\mathrm{vent}} \pi d_{\mathrm{vent}}.
\end{equation}
Using the same horizontal free-convection relation adopted for the open-top
case, the vent-opening convection coefficient is written as
\begin{equation}
    h_{\mathrm{top,vent}} = \frac{Nu\,k}{L_{c,\mathrm{vent}}}.
\end{equation}

With this mixed covered-open boundary condition, the top-node ambient-loss
conductance becomes
\begin{equation}
    UA_{\mathrm{top,covered+vent}} =
    U_{\mathrm{wall}}\left(\frac{1}{2}A_{\mathrm{side}} + A_{\mathrm{top,cov}}\right)
    + h_{\mathrm{top,vent}} A_{\mathrm{vent}}.
\end{equation}
The total vessel conductance to ambient is therefore
\begin{equation}
    (UA)_{\mathrm{covered+vent}} =
    U_{\mathrm{wall}}\left(A_{\mathrm{side}} + A_{\mathrm{bottom}} + A_{\mathrm{top,cov}}\right)
    + h_{\mathrm{top,vent}} A_{\mathrm{vent}}.
\end{equation}
The corresponding heat loss is
\begin{equation}
    Q_{\mathrm{loss,covered+vent}} =
    (UA)_{\mathrm{covered+vent}} \left(T_{\mathrm{bin}} - T_{\infty}\right).
\end{equation}

This formulation reduces to the fully covered-top model when
$A_{\mathrm{vent}} \rightarrow 0$, and it approaches the open-top model as
$A_{\mathrm{vent}} \rightarrow A_{\mathrm{top}}$. In the present
implementation, the vent is modeled as a thermal opening only; separate
effects such as moisture transport, buoyancy-driven exchange through the
vent, or forced vent exhaust are not added as independent loss terms.

\subsection{Perforation-Driven Moisture Transfer, Water Make-Up, and Computational Selection}

This section documents the moisture-transfer model, the lexicographic
optimization rule, and the numerical procedure used by the heater-design
code. The governing water-property relations follow the IAPWS saturation
ancillary equations, the hole-scale convection coefficient is obtained from
the Churchill--Bernstein crossflow correlation, and the heat--mass transfer
link is written with the Lewis relation. General energy balances and
continuity relations follow standard heat- and mass-transfer analysis.
In the source-code commentary these formula families are tagged as
IAPWS saturation relations [R9], Churchill--Bernstein crossflow [R1], the
Lewis relation [R10], standard energy/continuity balances [R5], and the
Frederickson et al.\ particle-size data [R11].

\subsubsection{Implemented Nusselt Correlations}

The code uses four distinct Nusselt-correlation branches, and the \LaTeX{}
documentation should state them explicitly because they are not interchangeable.

\paragraph*{1. Crossflow over the heater wire and perforation holes}

For both the heated wire and the perforation-scale crossflow, the code uses
the Churchill--Bernstein cylinder-in-crossflow relation [R1]:
\begin{equation}
Nu
=
0.3+
\left(
\frac{0.62\,Re^{1/2}Pr^{1/3}}
{\left[1+\left(0.4/Pr\right)^{2/3}\right]^{1/4}}
\right)
\left[1+\left(\frac{Re}{282000}\right)^{5/8}\right]^{4/5}.
\end{equation}
This is then converted to a convection coefficient by
\begin{equation}
h = \frac{Nu\,k}{D},
\end{equation}
with $D=D_{\mathrm{wire}}$ for wire cooling and $D=d_h$ for the
perforation-hole moisture calculation.

\paragraph*{2. Internal convection in the heater tube and aeration tubes}

The code uses a piecewise internal-flow model. For laminar developing flow,
it uses the Hausen-type entrance correction [R5]:
\begin{equation}
Gz = Re\,Pr\,\frac{D}{L},
\end{equation}
\begin{equation}
Nu
=
3.66+\frac{0.0668\,Gz}{1+0.04\,Gz^{2/3}},
\qquad Re<2300.
\end{equation}
For turbulent flow, it uses the Gnielinski correlation [R2]:
\begin{equation}
Nu
=
\frac{(f/8)(Re-1000)Pr}
{1+12.7(f/8)^{1/2}(Pr^{2/3}-1)},
\qquad Re\ge 2300.
\end{equation}
The implementation enforces
\begin{equation}
Nu \ge 3.66
\end{equation}
as a numerical lower bound. The same $Nu$ branch is used for both the heater
tube and the aeration tubes, and the same Churchill friction factor [R3]
supplies the $f$ entering the Gnielinski relation.

\paragraph*{3. Natural convection from the vessel vertical walls}

For the external vertical walls, the code evaluates the Rayleigh number
\begin{equation}
Ra = \frac{g\beta \Delta T L^3}{\nu \alpha},
\end{equation}
then uses the Churchill--Chu vertical-plate correlation [R4]:
\begin{equation}
Nu
=
\left[
0.825+
\frac{0.387\,Ra^{1/6}}
{\left(1+(0.492/Pr)^{9/16}\right)^{8/27}}
\right]^2.
\end{equation}
The resulting free-convection coefficient is
\begin{equation}
h = \frac{Nu\,k}{L}.
\end{equation}

\paragraph*{4. Natural convection from horizontal top or vent surfaces}

For the open top or vent opening, the code uses McAdams-style horizontal-plate
relations [R5]. For a heated upward-facing surface (the usual cooling/heating
case when the bed surface is warmer than ambient),
\begin{equation}
Nu = 0.54\,Ra^{1/4},
\qquad Ra \le 10^7,
\end{equation}
\begin{equation}
Nu = 0.15\,Ra^{1/3},
\qquad Ra > 10^7.
\end{equation}
For the opposite orientation branch used when the surface is cooler than the
surroundings, the code uses
\begin{equation}
Nu = 0.27\,Ra^{1/4}.
\end{equation}
Again,
\begin{equation}
h = \frac{Nu\,k}{L}.
\end{equation}

These four Nusselt branches are the complete set used by the present
implementation. Everything else in the thermal model uses those resulting
heat-transfer coefficients rather than introducing additional convection
correlations.

\subsubsection{Perforation Continuity and Hole Exit Velocity}

Let the total imposed inlet volumetric flow rate be $\dot{V}_{\mathrm{tot}}$,
measured at the inlet state. With $N_{\mathrm{tubes}}$ parallel aeration
tubes, the inlet flow per tube is
\begin{equation}
\dot{V}_{\mathrm{tube,in}} = \frac{\dot{V}_{\mathrm{tot}}}{N_{\mathrm{tubes}}}.
\end{equation}
Using the inlet air density $\rho_{\mathrm{in}}$, the inlet mass flow rate per
tube is
\begin{equation}
\dot{m}_{\mathrm{tube,in}} = \rho_{\mathrm{in}} \dot{V}_{\mathrm{tube,in}}.
\end{equation}

The model assumes that a fraction $\phi_{\mathrm{rel}}$ of the tube flow is
released into the bed through perforations, so that
\begin{equation}
\dot{m}_{\mathrm{rel}} = \phi_{\mathrm{rel}} \dot{m}_{\mathrm{tube,in}}.
\end{equation}
If each tube contains $N_h$ holes of diameter $d_h$, the area per hole is
\begin{equation}
A_h = \frac{\pi d_h^2}{4},
\end{equation}
and the total open area per tube is
\begin{equation}
A_{h,\mathrm{tot}} = N_h A_h.
\end{equation}

Because the released air is evaluated at the local heated-air state,
its volumetric release rate is
\begin{equation}
\dot{V}_{\mathrm{rel}} = \frac{\dot{m}_{\mathrm{rel}}}{\rho_{\mathrm{air}}(T_{\mathrm{air}},p)}.
\end{equation}
The mean released flow per hole is therefore
\begin{equation}
\dot{V}_{\mathrm{hole}} = \frac{\dot{V}_{\mathrm{rel}}}{N_h},
\end{equation}
and the mean hole exit velocity is
\begin{equation}
u_h = \frac{\dot{V}_{\mathrm{rel}}}{N_h A_h}.
\end{equation}

\paragraph*{Why $u_h$ varies with voltage.}
At fixed imposed inlet flow $\dot{V}_{\mathrm{tot}}$, the inlet mass flow rate
is nearly fixed by the inlet density. However, increasing the branch voltage
increases electrical power, which raises the heater outlet temperature
$T_{\mathrm{air}}$. Since the released-air density
$\rho_{\mathrm{air}}(T_{\mathrm{air}},p)$ decreases as the air is heated,
the same released mass flow occupies a larger volume. Consequently
$\dot{V}_{\mathrm{rel}}$ and hence $u_h$ increase with voltage even when the
imposed inlet flow axis is held constant.

\subsubsection{Particle-Size-Based Wetted-Area Estimate}

The current implementation no longer uses the bed plan area as the default
evaporation area. Instead, it derives an external particle-surface area from
the particle-size distribution reported in Table~3 of Frederickson et al.
[R11]. The reported matured-material size fractions are:
\begin{align}
\text{windrow: } &
9.8\% \ (>40\ \mathrm{mm}),\ 
26.6\% \ (20\text{--}40\ \mathrm{mm}),\ 
26.6\% \ (10\text{--}20\ \mathrm{mm}),\ 
37.0\% \ (<10\ \mathrm{mm}),\\
\text{vermicompost: } &
7.9\% \ (>40\ \mathrm{mm}),\ 
6.9\% \ (20\text{--}40\ \mathrm{mm}),\ 
19.8\% \ (10\text{--}20\ \mathrm{mm}),\ 
65.3\% \ (<10\ \mathrm{mm}).
\end{align}

The code forms the blended fractions
\begin{equation}
w_i = \frac{1}{2}\left(w_i^{\mathrm{windrow}} + w_i^{\mathrm{vermi}}\right),
\end{equation}
then assigns each bin an equivalent representative particle diameter
\begin{equation}
d_i \in \{60,\ 30,\ 15,\ 5\}\ \mathrm{mm}
\end{equation}
for the classes $>40$, $20$--$40$, $10$--$20$, and $<10$ mm respectively.
The closed bins use interval midpoints, while the open-ended $>40$ mm bin is
represented by $60$ mm, i.e.\ one adjacent-bin width above the lower cut.

The area calculation then uses the Sauter mean diameter for equal-density
equivalent spheres:
\begin{equation}
\frac{1}{d_{32}} = \sum_i \frac{w_i}{d_i},
\end{equation}
so that the external surface area per occupied bed volume is
\begin{equation}
a_{\mathrm{ext}} = \frac{6}{d_{32}}.
\end{equation}
This follows directly from the sphere area-to-volume ratio $A/V=6/d$ and is
the simplest particle-size-based estimate that still responds correctly to the
observed shift toward finer material.

Let the filled bed volume be
\begin{equation}
V_{\mathrm{fill}} = LWH f_{\mathrm{fill}}.
\end{equation}
The effective wetted area is then
\begin{equation}
A_{\mathrm{wet}} = f_{\mathrm{access}}\,a_{\mathrm{ext}}\,V_{\mathrm{fill}}
= f_{\mathrm{access}}\,\frac{6V_{\mathrm{fill}}}{d_{32}},
\end{equation}
where $f_{\mathrm{access}} \in [0,1]$ is the accessible-surface fraction used
to represent the share of the external particle area that is simultaneously
wet and exposed to the flowing air.

For the current default assumptions,
\begin{equation}
w_i = \{0.0885,\ 0.1675,\ 0.2320,\ 0.5115\},
\end{equation}
which gives
\begin{equation}
d_{32} \approx 8.01\ \mathrm{mm},
\qquad
a_{\mathrm{ext}} \approx 749\ \mathrm{m^2/m^3}.
\end{equation}
With
\begin{equation}
V_{\mathrm{fill}} = 1.22 \times 0.61 \times 0.61 \times 0.80
\approx 0.363\ \mathrm{m^3},
\end{equation}
and $f_{\mathrm{access}}=1$, the default effective wetted area becomes
\begin{equation}
A_{\mathrm{wet}} \approx 272\ \mathrm{m^2}.
\end{equation}

This is intentionally an \emph{external agglomerate surface} estimate, not a
true microscopic pore-wall area. The equal-density assumption means the
reported mass fractions are treated as proxies for volume fractions. That is
appropriate for a simple design-stage model because it preserves the dominant
effect of particle refinement on interfacial area without claiming a full
microstructural reconstruction of the compost matrix.

The code distributes this area uniformly across the aeration segments, so the
effective wetted area attached to a single segment is
\begin{equation}
A_{\mathrm{seg}}
= \frac{A_{\mathrm{wet}}}{N_{\mathrm{tubes}} N_{\mathrm{aer}}}.
\end{equation}

\subsubsection{Water Saturation Properties and Latent Heat}

The vapor-density driving force is built from saturation properties of water.
First, the saturation pressure is computed from the IAPWS ancillary form
\begin{equation}
\ln\!\left(\frac{p_{\mathrm{sat}}}{p_c}\right)
= \frac{T_c}{T}\sum_{i=1}^{6} a_i \tau^{b_i},
\qquad
\tau = 1 - \frac{T}{T_c},
\end{equation}
where $T_c$ and $p_c$ are the critical temperature and pressure of water.

The saturated liquid density is written as
\begin{equation}
\rho_{\ell,\mathrm{sat}}
= \rho_c \left(1 + \sum_{i=1}^{6} \beta_i \tau^{n_i}\right),
\end{equation}
and the saturated vapor density is
\begin{equation}
\rho_{v,\mathrm{sat}}
= \rho_c \exp\!\left(\sum_{i=1}^{6} \gamma_i \tau^{m_i}\right),
\end{equation}
with $\rho_c$ the critical density of water.

The saturation vapor density used in the mass-transfer driving force is then
\begin{equation}
\rho_{v,\mathrm{sat}}(T) = \frac{p_{\mathrm{sat}}(T)}{R_v T},
\end{equation}
where $R_v$ is the specific gas constant of water vapor.

The latent heat of vaporization is obtained from the Clausius--Clapeyron
relation along the saturation line:
\begin{equation}
h_{fg}(T)
= T \frac{dp_{\mathrm{sat}}}{dT}
\left(\frac{1}{\rho_{v,\mathrm{sat}}(T)} -
\frac{1}{\rho_{\ell,\mathrm{sat}}(T)}\right).
\end{equation}

\subsubsection{Hole-Scale Convection and Mass-Transfer Coefficients}

The model treats each perforation as a characteristic crossflow scale for
heat and mass transfer. Using the local heated-air properties, the hole-scale
Reynolds number is
\begin{equation}
\mathrm{Re}_h
= \frac{\rho_{\mathrm{air}} u_h d_h}{\mu_{\mathrm{air}}},
\end{equation}
and the Nusselt number is evaluated with the Churchill--Bernstein correlation:
\begin{equation}
\mathrm{Nu}_h
= 0.3 +
\frac{0.62 \mathrm{Re}_h^{1/2} \mathrm{Pr}^{1/3}}
{\left[1 + (0.4/\mathrm{Pr})^{2/3}\right]^{1/4}}
\left[1 + \left(\frac{\mathrm{Re}_h}{282000}\right)^{5/8}\right]^{4/5}.
\end{equation}
The corresponding convective coefficient is
\begin{equation}
h_{\mathrm{evap}}
= \mathrm{Nu}_h \frac{k_{\mathrm{air}}}{d_h}.
\end{equation}

The heat--mass transfer analogy is then used in the form
\begin{equation}
h_m
= \frac{h_{\mathrm{evap}}}
{\rho_{\mathrm{air}} c_{p,\mathrm{air}} \mathrm{Le}^{2/3}},
\end{equation}
where $\mathrm{Le}$ is the Lewis factor. In the present implementation,
$\mathrm{Le}$ is taken near unity for the air--water system.

\subsubsection{Evaporation Driving Force and Three Limiting Rates}

Let the prescribed effective wet-surface temperature be $T_s$ and the inlet
relative humidity be $\phi_{\infty}$. The model constructs the vapor-density
driving force from
\begin{equation}
\rho_{v,s} = \rho_{v,\mathrm{sat}}(T_s),
\end{equation}
\begin{equation}
\rho_{v,\infty}
= \phi_{\infty} \rho_{v,\mathrm{sat}}(T_{\infty}),
\end{equation}
so that
\begin{equation}
\Delta \rho_v
= \max\!\left(\rho_{v,s} - \rho_{v,\infty}, 0\right).
\end{equation}

For each aeration-tube segment, the model computes three candidate
evaporation rates and accepts the smallest of the three.

\paragraph*{Diffusion-limited rate.}
If $A_{\mathrm{seg}}$ is the effective wetted area assigned to a single
segment, then
\begin{equation}
\dot{m}_{\mathrm{diff}}
= h_m A_{\mathrm{seg}} \Delta \rho_v.
\end{equation}

\paragraph*{Carrying-capacity limit of the released air parcel.}
If $\dot{V}_{\mathrm{rel,seg}}$ is the released volumetric flow rate for that
segment, then the maximum vapor uptake permitted by the released air volume is
\begin{equation}
\dot{m}_{\mathrm{cap}}
= \dot{V}_{\mathrm{rel,seg}} \Delta \rho_v.
\end{equation}

\paragraph*{Heat-limited evaporation rate.}
The released jet also has a finite sensible heat content. The model defines
the available sensible power for evaporation as
\begin{equation}
q_{\mathrm{avail}}
= \max\!\left(
\dot{m}_{\mathrm{rel,seg}} c_{p,\mathrm{air}}
(T_{\mathrm{air}} - T_{\mathrm{bed}}), 0
\right),
\end{equation}
which leads to
\begin{equation}
\dot{m}_{\mathrm{heat}}
= \frac{q_{\mathrm{avail}}}{h_{fg}}.
\end{equation}

The accepted segment evaporation rate is
\begin{equation}
\dot{m}_{\mathrm{evap,seg}}
= \min\!\left(
\dot{m}_{\mathrm{diff}},
\dot{m}_{\mathrm{cap}},
\dot{m}_{\mathrm{heat}}
\right),
\end{equation}
and the corresponding latent heat sink is
\begin{equation}
q_{\mathrm{evap,seg}}
= \dot{m}_{\mathrm{evap,seg}} h_{fg}.
\end{equation}

\paragraph*{Unit consistency.}
The three candidate evaporation rates are dimensionally consistent:
\begin{itemize}
\item $h_m A \Delta \rho_v$ has units
      $(\mathrm{m/s})(\mathrm{m^2})(\mathrm{kg/m^3})=\mathrm{kg/s}$;
\item $\dot{V}_{\mathrm{rel,seg}} \Delta \rho_v$ has units
      $(\mathrm{m^3/s})(\mathrm{kg/m^3})=\mathrm{kg/s}$; and
\item $q_{\mathrm{avail}}/h_{fg}$ has units
      $(\mathrm{J/s})/(\mathrm{J/kg})=\mathrm{kg/s}$.
\end{itemize}
Therefore $\dot{m}_{\mathrm{diff}}$, $\dot{m}_{\mathrm{cap}}$,
$\dot{m}_{\mathrm{heat}}$, and $\dot{m}_{\mathrm{evap,seg}}$ are all
compared in consistent $\mathrm{kg/s}$ units.

\paragraph*{Why water loss and latent load vary with voltage.}
Voltage increases electrical heating, which typically raises
$T_{\mathrm{air}}$ at fixed inlet flow. This has three effects:
\begin{enumerate}
\item It lowers the local air density, thereby increasing
      $\dot{V}_{\mathrm{rel}}$ and $u_h$.
\item It increases $\mathrm{Re}_h$ and hence usually increases
      $h_{\mathrm{evap}}$ and $h_m$.
\item It increases the sensible heating term $q_{\mathrm{avail}}$, which
      raises the heat-limited evaporation cap.
\end{enumerate}
Accordingly, the water-loss and latent-load maps generally increase with
voltage, although the final contours are shaped by the minimum operator among
the diffusion, carrying-capacity, and heat-limited branches.

\subsubsection{Radial Air-Cooling Profile and Worm-Safe Distance}

The standalone Python study now adds a radial post-processing model to
estimate how rapidly the released warm air cools as it moves outward from the
aeration tube into the wet bed. This is not a new CFD solver; it is a finite
top-hat plume model built from the same moisture and wetted-area quantities
already computed by the main design code.

Let $x$ denote the radial distance measured outward from the aeration-tube
outer surface, and let
\begin{equation}
r_0 = \frac{D_{\mathrm{tube,OD}}}{2}
\end{equation}
be the tube outer radius. The released mass flow rate per unit tube length is
\begin{equation}
\dot{m}'_{\mathrm{rel}}
= \frac{\phi_{\mathrm{rel}} \dot{m}_{\mathrm{tube,in}}}{L_{\mathrm{aer}}}.
\end{equation}

Using the previously derived wetted area, the corresponding interfacial area
density is
\begin{equation}
a_{\mathrm{wet}} = \frac{A_{\mathrm{wet}}}{V_{\mathrm{fill}}}.
\end{equation}
Here $T_{\mathrm{sink}}$ is the clamped wet-surface temperature used as the
local cooling asymptote:
\begin{equation}
T_{\mathrm{sink}}
= \min\!\left(\max(T_{s,\mathrm{set}},T_{\mathrm{bed,ref}}),T_{\mathrm{out}}\right).
\end{equation}

Instead of assuming that the full local wetted area is active immediately,
the plume model assumes the warm released air occupies a finite plume of
radius $b(x)$ that grows with distance:
\begin{equation}
b(x) = \min\!\left(b_{\max},\, b_0 + \alpha x\right).
\end{equation}
Here $b_0$ is the initial plume radius, usually taken as the hole radius,
$\alpha$ is the plume spreading rate, and $b_{\max}$ is an optional cap.
These are exposed in the JSON file through the
\texttt{radial\_profile} block.

Unlike the earlier approximation, the present implementation does
\emph{not} collapse latent cooling into a single fixed equivalent
coefficient. Instead it re-solves the plume temperature and vapor state
along the radial path. The contacted wetted area per unit radial distance is
\begin{equation}
A'_{\mathrm{wet}}(x) = \pi a_{\mathrm{wet}} b(x)^2.
\end{equation}
The sensible cooling term is
\begin{equation}
q'_{\mathrm{sens}}
= h_{\mathrm{sens}} A'_{\mathrm{wet}}(x)\left(T-T_{\mathrm{sink}}\right),
\end{equation}
and the temperature equation is
\begin{equation}
\dot{m}'_{\mathrm{rel}} c_{p,\mathrm{air}} \frac{dT}{dx}
= -\left(q'_{\mathrm{sens}} + q'_{\mathrm{lat}}\right).
\end{equation}

The vapor-density state is updated from a Lewis-relation mass-transfer model
written with local air properties:
\begin{equation}
h_m(x)
=
\frac{h_{\mathrm{evap}}}
{\rho_{\mathrm{air}}(T)\,c_{p,\mathrm{air}}\,\mathrm{Le}^{2/3}}.
\end{equation}
Using the same wet-surface temperature clamp as the main moisture model, the
local diffusion-driven evaporation rate per unit radial distance is
\begin{equation}
\dot{m}'_{\mathrm{evap,pot}}
=
h_m(x)\,A'_{\mathrm{wet}}(x)\,
\left(\rho_{v,\mathrm{sat}}(T_s)-\rho_{v,\mathrm{bulk}}(x)\right).
\end{equation}
The plume humidity is then advanced explicitly so that the bulk vapor density
cannot exceed local saturation at the updated plume temperature. In this
way, the latent sink weakens automatically as the plume both cools and
humidifies.

The corresponding latent heat term is
\begin{equation}
q'_{\mathrm{lat}}
=
\dot{m}'_{\mathrm{evap}} h_{fg}(T_s),
\end{equation}
and the total temperature drop is limited so that the plume cannot cool
below $T_{\mathrm{sink}}$ within a single marching step.

This is still a reduced-order model rather than CFD, but it is more
physically consistent than the earlier constant-$h_{\mathrm{eff}}$
approximation because the radial evaporation rate now responds directly to
the local plume temperature and local vapor build-up.
\par
The solution is obtained numerically on the radial grid used for the plot;
the earlier closed-form expression for $T(x)$ is no longer used because the
latent term is now state-dependent.
\par
If the worm-safety threshold is $T_{\mathrm{safe}} = 27^\circ\mathrm{C}$,
the safe offset from the tube outer surface is the first $x$ satisfying
\begin{equation}
T(x)=T_{\mathrm{safe}}.
\end{equation}
In the code this crossing is found by interpolation on the numerically
evaluated profile.

This construction is intentionally modest in scope. It still does
\emph{not} resolve individual jets or local tortuous pore paths, but it is
more physically consistent than the earlier immediate-full-area sink because
it delays plume coupling to the full local wetted surface until the warm-air
core has spread through a larger pore volume.

\subsubsection{Tube Layout, Safe-Distance Exclusion Cylinders, and Habitable Volume}

The standalone study also converts the worm-safe radial cutoff into a
geometric estimate of the bed volume that is effectively too warm for worm
occupation. This calculation is performed in the bed cross section and then
extruded over the aeration-tube length.

Let the filled-bed height be
\begin{equation}
H_{\mathrm{fill}} = f_{\mathrm{fill}} H,
\end{equation}
so that the filled-bed cross-sectional domain is
\begin{equation}
\Omega_{\mathrm{fill}}
=
\left\{(x,y): 0 \le x \le W,\ 0 \le y \le H_{\mathrm{fill}}\right\}.
\end{equation}
The corresponding filled-bed volume is
\begin{equation}
V_{\mathrm{fill}} = L W H_{\mathrm{fill}}.
\end{equation}

If the user does not specify tube-center locations explicitly, the current
implementation uses the requested rule-based layout for $1 \le
N_{\mathrm{tubes}} \le 8$ inside the \emph{filled-bed} rectangle. Let
\begin{equation}
x_L = f_x W, \qquad x_R = (1-f_x)W,
\end{equation}
\begin{equation}
y_B = f_y H_{\mathrm{fill}}, \qquad y_T = (1-f_y)H_{\mathrm{fill}},
\end{equation}
with default offsets $f_x=f_y=0.10$. The rules are:
\begin{itemize}
\item 1 tube: $(W/2,\ y_B)$;
\item 2 tubes: bottom-left and bottom-right;
\item 3 tubes: three equally spaced across the bottom row from $x_L$ to $x_R$;
\item 4 tubes: the four corners $(x_L,y_B)$, $(x_R,y_B)$, $(x_L,y_T)$, $(x_R,y_T)$;
\item 5 tubes: the 3-bottom-row pattern plus the two top corners;
\item 6 tubes: three equally spaced vertically on the left side and three on the right side;
\item 7 tubes: the four corners plus left-mid, right-mid, and bottom-center; and
\item 8 tubes: the 7-tube pattern plus a central tube at $(W/2,\ H_{\mathrm{fill}}/2)$.
\end{itemize}
If explicit positions are supplied in the JSON file, those values replace the
rule-based pattern directly.

The center-to-center pitch between neighboring tubes is
\begin{equation}
p_i = x_{i+1} - x_i,
\end{equation}
and the corresponding clear gap between the outer tube surfaces is
\begin{equation}
g_i = p_i - D_{\mathrm{tube,OD}}.
\end{equation}

Using the radial safe offset $x_{\mathrm{safe}}$ from the previous section,
the safe-distance exclusion radius around each tube is
\begin{equation}
R_{\mathrm{exc}} = r_0 + x_{\mathrm{safe}}.
\end{equation}
Thus each actual pipe occupies a circular cross section of radius $r_0$,
while the worm-exclusion region occupies a larger circle of radius
$R_{\mathrm{exc}}$ centered on the same tube axis.

To account correctly for overlap between neighboring exclusion envelopes, the
study does \emph{not} sum four independent cylinders. Instead, it evaluates
the union area in cross section. Define the pipe-indicator function
\begin{equation}
I_{\mathrm{pipe}}(x,y)
=
\begin{cases}
1, & \exists i \text{ such that } (x-x_i)^2 + (y-y_i)^2 \le r_0^2, \\
0, & \text{otherwise},
\end{cases}
\end{equation}
and the exclusion-indicator function
\begin{equation}
I_{\mathrm{exc}}(x,y)
=
\begin{cases}
1, & \exists i \text{ such that } (x-x_i)^2 + (y-y_i)^2 \le R_{\mathrm{exc}}^2, \\
0, & \text{otherwise}.
\end{cases}
\end{equation}
The corresponding overlap-aware union areas are
\begin{equation}
A_{\mathrm{pipe,union}}
=
\iint_{\Omega_{\mathrm{fill}}} I_{\mathrm{pipe}}(x,y)\, dA,
\end{equation}
and
\begin{equation}
A_{\mathrm{exc,union}}
=
\iint_{\Omega_{\mathrm{fill}}} I_{\mathrm{exc}}(x,y)\, dA.
\end{equation}

Since the aeration tubes run parallel to the vessel length, the associated
volumes are obtained by extrusion:
\begin{equation}
V_{\mathrm{pipe}} = A_{\mathrm{pipe,union}} L_{\mathrm{aer}},
\end{equation}
\begin{equation}
V_{\mathrm{exc}} = A_{\mathrm{exc,union}} L_{\mathrm{aer}}.
\end{equation}
The pipe-solid volume is therefore $V_{\mathrm{pipe}}$, the worm-unsafe shell
volume outside the pipe metal but inside the safe-distance envelope is
\begin{equation}
V_{\mathrm{unsafe,shell}} = V_{\mathrm{exc}} - V_{\mathrm{pipe}},
\end{equation}
and the remaining habitable bed volume is
\begin{equation}
V_{\mathrm{hab}} = V_{\mathrm{fill}} - V_{\mathrm{exc}}.
\end{equation}

The current implementation also adds a simple wall-return correction for the
left, right, and bottom vessel sheets. If the plume reaches a wall, the wall
contact temperature is sampled from the radial plume solution at the
wall-offset distance. Using the wall stack already present in the main heat
loss model, the inside wall temperature is estimated from the 1D resistance
network
\begin{equation}
q''_{\mathrm{wall}}
=
\frac{T_{\mathrm{plume,wall}} - T_{\infty}}
{\frac{1}{h_{\mathrm{in}}}
+ \frac{t_{\mathrm{sheet}}}{k_{\mathrm{sheet}}}
+ \frac{t_{\mathrm{ins}}}{k_{\mathrm{ins}}}
+ \frac{1}{h_{\mathrm{ext}}}},
\end{equation}
\begin{equation}
T_{\mathrm{wall,in}}
=
T_{\mathrm{plume,wall}} - \frac{q''_{\mathrm{wall}}}{h_{\mathrm{in}}}.
\end{equation}
Assuming the steel sheet redistributes this temperature approximately
uniformly along the corresponding wall, a 1D linear wall-normal soil profile
is then used to estimate an additional unsafe band depth. Those wall bands
are merged with the exclusion cylinders through the same union-area
calculation, so the final unsafe volume reported by the code is the union of
pipe cylinders, safe-distance cylinders, and wall-return bands:
\begin{equation}
V_{\mathrm{unsafe,total}} = V_{\mathrm{union}}.
\end{equation}

The code also reports explicit overlap corrections by comparing the union
areas above with the naive no-overlap sums:
\begin{equation}
A_{\mathrm{pipe,sum}} = N_{\mathrm{tubes}} \pi r_0^2,
\qquad
A_{\mathrm{exc,sum}} = N_{\mathrm{tubes}} \pi R_{\mathrm{exc}}^2.
\end{equation}
The removed double-counted areas are
\begin{equation}
A_{\mathrm{pipe,ov}} = A_{\mathrm{pipe,sum}} - A_{\mathrm{pipe,union}},
\end{equation}
\begin{equation}
A_{\mathrm{exc,ov}} = A_{\mathrm{exc,sum}} - A_{\mathrm{exc,union}},
\end{equation}
with volume counterparts
\begin{equation}
V_{\mathrm{pipe,ov}} = A_{\mathrm{pipe,ov}} L_{\mathrm{aer}},
\qquad
V_{\mathrm{exc,ov}} = A_{\mathrm{exc,ov}} L_{\mathrm{aer}}.
\end{equation}
These overlap terms are useful because they show how much the naive
``sum of independent cylinders'' estimate would over-predict the inaccessible or unsafe
volume when neighboring exclusion envelopes intersect.

Numerically, the present implementation evaluates
$A_{\mathrm{pipe,union}}$ and $A_{\mathrm{exc,union}}$ on a midpoint grid
over $\Omega_{\mathrm{fill}}$. If $N_g$ points are used in each direction,
the area estimator is
\begin{equation}
A \approx \left(\frac{W H_{\mathrm{fill}}}{N_g^2}\right)
\sum_{j=1}^{N_g} \sum_{k=1}^{N_g} I(x_j^\ast,y_k^\ast),
\end{equation}
where $(x_j^\ast,y_k^\ast)$ are the cell midpoints. The resulting union area
is then multiplied by $L_{\mathrm{aer}}$ to obtain the reported volume.
This is sufficient here because the cross section contains only a small
number of circles and the output is intended as an engineering habitat-volume
estimate rather than an exact computational-geometry benchmark.

\subsubsection{Aeration-Tube Energy Balance with Moisture Sink}

For each aeration-tube segment, the wall-mediated sensible transfer is
\begin{equation}
q_{\mathrm{wall,seg}}
= U_{\mathrm{aer}} P_{\mathrm{out}} \Delta x
(T_{\mathrm{air}} - T_{\mathrm{bed}}),
\end{equation}
and the direct sensible jet contribution is
\begin{equation}
q_{\mathrm{jet,seg}}
= \dot{m}_{\mathrm{rel,seg}} c_{p,\mathrm{air}}
(T_{\mathrm{air}} - T_{\mathrm{bed}}).
\end{equation}
With evaporation included explicitly, the useful heat delivered to the bed by
one aeration branch is
\begin{equation}
Q_{\mathrm{bed,branch}}
= \sum_{\mathrm{seg}}
\left(
q_{\mathrm{wall,seg}} +
q_{\mathrm{jet,seg}} -
q_{\mathrm{evap,seg}}
\right).
\end{equation}
For $N_{\mathrm{tubes}}$ identical parallel tubes,
\begin{equation}
Q_{\mathrm{bed,total}}
= N_{\mathrm{tubes}} Q_{\mathrm{bed,branch}}.
\end{equation}

The total modeled water make-up requirement is
\begin{equation}
\dot{m}_{\mathrm{water,day}}
= 86400\,
N_{\mathrm{tubes}}
\sum_{\mathrm{seg}}
\dot{m}_{\mathrm{evap,seg}},
\end{equation}
reported in $\mathrm{kg/24\,h}$. This quantity is treated as the sprinkler
replacement requirement and is also used directly as a design constraint.

\subsection{Lexicographic Optimization and Design-Point Selection}

The code does not use a weighted penalty function. Instead, every candidate
operating point is assigned an ordered criteria vector and candidate points
are compared lexicographically. Thus the first unequal component decides the
ordering, and later criteria cannot compensate for an earlier violation.

\subsubsection{Candidate Criteria Vector}

For each operating point, the code first computes the full-power two-node bed
temperatures $T_b$ and $T_t$, their spread
\begin{equation}
\Delta T_{\mathrm{bed}} = |T_b - T_t|,
\end{equation}
and the duty ratio required to meet the lumped design heat loss,
\begin{equation}
\mathrm{dutyNeeded}
= \frac{Q_{\mathrm{req}}}{Q_{\mathrm{bed,total}}}.
\end{equation}

It then constructs the ordered criteria vector
\begin{equation}
\mathbf{C}
=
\left[
C_1,C_2,C_3,C_4,C_5,C_6,C_7,C_8,C_9,C_{10},C_{11},C_{12},C_{13}
\right],
\end{equation}
with
\begin{align}
C_1 &= \max(I_{\mathrm{tot}} - I_{\max}, 0), \\
C_2 &= \max(T_{\mathrm{wire}} - T_{\mathrm{wire,max}}, 0), \\
C_3 &= \max(T_{\mathrm{air,out}} - T_{\mathrm{air,max}}, 0), \\
C_4 &= \max(\Delta P - \Delta P_{\max}, 0), \\
C_5 &= \max(\dot{m}_{\mathrm{water,day}} - \dot{m}_{\mathrm{water,max}}, 0), \\
C_6 &= \max(T_{b,\min} - T_b, 0), \\
C_7 &= \max(T_{t,\min} - T_t, 0), \\
C_8 &= \max(\Delta T_{\mathrm{bed}} - \Delta T_{\max}, 0), \\
C_9 &= -\min(T_b,T_t), \\
C_{10} &= P_{\mathrm{tot}}, \\
C_{11} &= \Delta T_{\mathrm{bed}}, \\
C_{12} &= -\frac{T_b + T_t}{2}, \\
C_{13} &= \max(\mathrm{dutyNeeded} - 1, 0).
\end{align}

The first eight entries are hard-constraint residuals. If any of them is
positive, the point is not hard-safe. Among hard-safe points, the clean
selection rule is therefore:
\begin{enumerate}
\item maximize the colder-node temperature $\min(T_b,T_t)$;
\item minimize total electrical power $P_{\mathrm{tot}}$;
\item minimize bottom--top spread $\Delta T_{\mathrm{bed}}$; and
\item use mean bed temperature and residual heat shortfall only as late
      tie-breakers.
\end{enumerate}
This ordering is more physically transparent than a weighted score because it
states directly which tradeoffs are acceptable and in what order.

\subsubsection{Formal Lexicographic Rule}

Given two candidate points $A$ and $B$ with criteria vectors
$\mathbf{C}^{(A)}$ and $\mathbf{C}^{(B)}$, point $A$ is preferred if there
exists an index $k$ such that
\begin{equation}
C^{(A)}_j = C^{(B)}_j \quad \text{for all } j < k,
\end{equation}
and
\begin{equation}
C^{(A)}_k < C^{(B)}_k.
\end{equation}
This is exactly the row-wise ordering used by the implementation.

\subsubsection{Feasible, Fallback, and Last-Resort Classifications}

The code additionally labels each point by its logical stage:
\begin{enumerate}
\item \textbf{Feasible:} all explicit constraints are satisfied and
      $\mathrm{dutyNeeded} \le 1$.
\item \textbf{Fallback:} all hard constraints are satisfied and
      $Q_{\mathrm{bed,total}} > 0$, but the full heating target is not met.
\item \textbf{Last resort:} one or more explicit constraints are violated,
      or the net delivered heat is non-positive.
\end{enumerate}
These labels are descriptive only. The actual selection still uses the full
lexicographic criteria vector above. Thus the recommended point is always the
best hard-safe point if one exists, even when no fully feasible point exists.
Separately, the report also carries a configurable duty-reference value
\texttt{model\_overrides.sweep.targetDuty}. This quantity does \emph{not}
change the hard-feasibility definition above; instead it controls the
reported line ``minimum total flow at \ldots duty'' and the summary note about
whether the selected point can meet that user-specified partial-duty target.
In the current configuration this report-only duty reference is set to
$0.75$, i.e. $75\%$ duty.
The editable geometry/material blocks that change the underlying exported
physics are \texttt{model\_overrides.bin\_wall},
\texttt{model\_overrides.heater\_tube}, \texttt{model\_overrides.aeration},
and \texttt{model\_overrides.heaterTube}; the corresponding
\texttt{summary\_inputs} entries are fallback report text only.

\subsubsection{Where the Same Rule Is Used}

The exact same comparison rule is applied in four places:
\begin{enumerate}
\item selection of the reported recommended point for each vessel
      configuration;
\item ranking of the top candidate points for the active configuration;
\item selection of the single representative point at each sampled wire
      diameter in the wire-diameter study; and
\item generation of the optimization-ranked table.
\end{enumerate}
Hence the report contains one internally consistent ordering principle rather
than separate weighted objectives for different figures.

\subsection{Numerical Procedure}

\subsubsection{Axial Discretization of the Heater and Aeration Tubes}

The heater tube is discretized into $N_{\mathrm{heater}}$ axial segments,
with
\begin{equation}
\Delta x_{\mathrm{heater}}
= \frac{L_{\mathrm{heater}}}{N_{\mathrm{heater}}},
\end{equation}
and the aeration tube is discretized into $N_{\mathrm{aer}}$ axial segments,
with
\begin{equation}
\Delta x_{\mathrm{aer}}
= \frac{L_{\mathrm{aer}}}{N_{\mathrm{aer}}}.
\end{equation}
In the current implementation,
\begin{equation}
N_{\mathrm{heater}} = 50,
\qquad
N_{\mathrm{aer}} = 40.
\end{equation}
The solution marches explicitly from inlet to outlet through both domains.

\subsubsection{Fixed-Point Iteration for Temperature-Dependent Wire Resistance}

For a specified branch voltage $V$ and total flow $\dot{V}_{\mathrm{tot}}$,
the code solves the operating point by iterating on the average wire
temperature. Starting from an initial guess
$\overline{T}_{\mathrm{wire}}^{(0)}$, it evaluates the temperature-dependent
wire resistance
\begin{equation}
\rho_{\mathrm{wire}}(T)
= \rho_{\mathrm{wire},20}
\left[1 + \alpha_{\mathrm{wire}}(T-20^\circ\mathrm{C})\right],
\end{equation}
\begin{equation}
R_{\mathrm{branch}}
= \rho_{\mathrm{wire}}(T)
\frac{L_{\mathrm{wire}}}{A_{\mathrm{wire}}},
\end{equation}
then computes
\begin{equation}
I_{\mathrm{branch}} = \frac{V}{R_{\mathrm{branch}}},
\qquad
P_{\mathrm{branch}} = V I_{\mathrm{branch}}.
\end{equation}

The heater-tube solution then returns a new mean wire temperature
$\overline{T}_{\mathrm{wire,new}}$. The code updates the iterate by damping:
\begin{equation}
\overline{T}_{\mathrm{wire}}^{(k+1)}
= 0.65\,\overline{T}_{\mathrm{wire}}^{(k)}
+ 0.35\,\overline{T}_{\mathrm{wire,new}}^{(k)}.
\end{equation}
Convergence is accepted when
\begin{equation}
\left|
\overline{T}_{\mathrm{wire,new}}^{(k)} -
\overline{T}_{\mathrm{wire}}^{(k)}
\right| < 0.05^\circ\mathrm{C},
\end{equation}
or when the maximum iteration count is reached. The present implementation
uses at most eight iterations per operating point.

The coefficients $0.65$ and $0.35$ are numerical under-relaxation factors,
not physical material parameters. They are used because the coupled
resistance--current--wire-temperature loop can oscillate if the full updated
temperature is accepted at every iteration. A hotter guessed wire raises the
resistance, which lowers current and power, which then lowers the predicted
wire temperature on the next pass. Under-relaxation damps this feedback by
blending the old iterate with the newly computed iterate. The chosen values
are empirical but moderate: they preserve convergence speed while suppressing
oscillatory overshoot over the voltage-flow sweep.

\subsubsection{Heater-Tube Marching Scheme}

Within each heater segment, the local generated electrical heat is
\begin{equation}
q_{\mathrm{gen,seg}}
= P_{\mathrm{branch}}
\frac{\Delta x_{\mathrm{heater}}}{L_{\mathrm{heater}}},
\end{equation}
the heater-shell loss to ambient is
\begin{equation}
q_{\mathrm{loss,seg}}
= U_{\mathrm{heater}} P_{\mathrm{out}} \Delta x_{\mathrm{heater}}
(T_{\mathrm{air}} - T_{\infty}),
\end{equation}
and the net heat retained by the airflow is
\begin{equation}
q_{\mathrm{to\,air,seg}}
= q_{\mathrm{gen,seg}} - q_{\mathrm{loss,seg}}.
\end{equation}
The air temperature is advanced explicitly by
\begin{equation}
T_{\mathrm{air},i+1}
= T_{\mathrm{air},i}
+ \frac{q_{\mathrm{to\,air,seg}}}
{\dot{m}_{\mathrm{heater}} c_{p,\mathrm{air}}}.
\end{equation}
The local wire temperature is then reconstructed from Newton's law:
\begin{equation}
T_{\mathrm{wire},i}
= T_{\mathrm{air},i}
+ \frac{q_{\mathrm{gen,seg}}}
{h_{\mathrm{wire}} A_{\mathrm{wire,seg}}}.
\end{equation}

Here $T_{\infty}$ denotes the surrounding greenhouse ambient used by the
heater-shell loss model. In the present study this is the design greenhouse
air temperature (for example $-15^\circ\mathrm{C}$), not a separately solved
local external air temperature field around the heater tube.

\subsubsection{Aeration-Tube Marching Scheme with Progressive Release}

The aeration tube receives the heater outlet state as its inlet. The total
released mass is distributed uniformly over the aeration segments:
\begin{equation}
\dot{m}_{\mathrm{rel,seg}}
= \frac{\phi_{\mathrm{rel}} \dot{m}_{\mathrm{tube,in}}}
{N_{\mathrm{aer}}}.
\end{equation}
At each segment, the code computes
\begin{equation}
q_{\mathrm{wall,seg}}
= U_{\mathrm{aer}} P_{\mathrm{out}} \Delta x_{\mathrm{aer}}
(T_{\mathrm{air}} - T_{\mathrm{bed}}),
\end{equation}
\begin{equation}
q_{\mathrm{jet,seg}}
= \dot{m}_{\mathrm{rel,seg}} c_{p,\mathrm{air}}
(T_{\mathrm{air}} - T_{\mathrm{bed}}),
\end{equation}
\begin{equation}
q_{\mathrm{evap,seg}}
= \dot{m}_{\mathrm{evap,seg}} h_{fg},
\end{equation}
and accumulates
\begin{equation}
Q_{\mathrm{bed,branch}}
\leftarrow
Q_{\mathrm{bed,branch}}
+ q_{\mathrm{wall,seg}} + q_{\mathrm{jet,seg}} - q_{\mathrm{evap,seg}}.
\end{equation}

The remaining in-tube mass flow is reduced after each segment:
\begin{equation}
\dot{m}_{\mathrm{rem},i+1}
= \dot{m}_{\mathrm{rem},i} - \dot{m}_{\mathrm{rel,seg}},
\end{equation}
and the bulk-tube air temperature is advanced using the wall-transfer term:
\begin{equation}
T_{\mathrm{air},i+1}
= T_{\mathrm{air},i}
- \frac{q_{\mathrm{wall,seg}}}
{\dot{m}_{\mathrm{rem},i+1} c_{p,\mathrm{air}}}.
\end{equation}

\subsubsection{Pressure-Drop Evaluation}

Both the heater and aeration tubes accumulate Darcy--Weisbach friction and
distributed minor losses segment by segment:
\begin{equation}
\Delta P_{\mathrm{seg}}
= \left(
f \frac{\Delta x}{D} + \frac{K_{\mathrm{minor}}}{N}
\right)
\frac{\rho u^2}{2}.
\end{equation}
The heater-to-aeration abrupt expansion is modeled by
\begin{equation}
\Delta P_{\mathrm{exp}}
= \frac{\rho}{2}(u_1-u_2)^2,
\end{equation}
which is the standard Borda--Carnot sudden-expansion relation.

For the common distribution hardware upstream of each branch, the current
implementation distinguishes five contributions:
\begin{enumerate}
\item an optional upstream header loss,
\item an optional header-to-splitter sudden contraction,
\item a splitter-body minor loss,
\item friction in the short branch connector, and
\item a connector-to-branch sudden contraction or sudden expansion.
\end{enumerate}

If a separate large upstream header is used, the header term is
\begin{equation}
\Delta P_{\mathrm{header}}
= K_{\mathrm{header}} \frac{\rho u_{\mathrm{header}}^2}{2}.
\end{equation}
If \texttt{headerEnabled = false}, this term is omitted.

Let the nominal splitter/connector diameter be $D_c$, let the branch
downstream diameter be $D_b$, and let the short connector length be
$L_c$. The connector-friction term is
\begin{equation}
\Delta P_{\mathrm{connector}}
= f_c \frac{L_c}{D_c}\frac{\rho u_c^2}{2},
\end{equation}
with $u_c$ and $f_c$ computed from the connector hydraulic diameter and the
branch mass flow rate.

For a conical converging bellmouth without end wall, as tabulated by
Idel'chik in Section~III, Diagram~3-5, the code models a contraction from
an upstream diameter $D_1$ to a downstream diameter $D_2 < D_1$ by
defining
\begin{equation}
\epsilon = \frac{A_2}{A_1},
\end{equation}
reading the diagram coefficient $C(\epsilon,\alpha)$ at the configured cone
angle $\alpha$, and then using Idel'chik's equation~(3-3) on the smaller-pipe
velocity basis:
\begin{equation}
\Delta P_{\mathrm{contr}}
= K_{\mathrm{contr}} \frac{\rho u_2^2}{2},
\end{equation}
with
\begin{equation}
\alpha = 60^\circ \quad \text{(current default)},
\qquad
K_{\mathrm{contr}} = C(\epsilon,\alpha)\left(\frac{1}{\epsilon}-1\right).
\end{equation}
The current implementation interpolates the tabulated $C$ values from
Diagram~3-5 over
\begin{equation}
\epsilon \in \{0.025, 0.050, 0.075, 0.10, 0.15, 0.25, 0.60, 1.0\}
\end{equation}
and
\begin{equation}
\alpha \in \{0^\circ, 10^\circ, 20^\circ, 30^\circ, 40^\circ, 60^\circ,
100^\circ, 140^\circ, 180^\circ\}.
\end{equation}
This form is used for the optional header-to-splitter contraction and, when
applicable, for the connector-to-downstream contraction.

For a sudden expansion from $D_1$ to $D_2 > D_1$, the code uses the
Idel'chik Section~IV sudden-expansion coefficient, written on the
upstream-velocity basis as
\begin{equation}
\Delta P_{\mathrm{exp}}
= K_{\mathrm{exp}} \frac{\rho u_1^2}{2},
\qquad
K_{\mathrm{exp}} = \left(1-\frac{A_1}{A_2}\right)^2,
\end{equation}
which is algebraically equivalent to the Borda--Carnot form already written
above and is the same coefficient quoted in standard texts such as White and
Munson when the loss is referenced to the upstream velocity.

The internal loss of the commercial $4$-way splitter body is not published
by the vendor. Therefore the code represents that unresolved internal
branching loss by the editable lumped coefficient
\texttt{splitterBodyK}, applied on the connector-velocity basis:
\begin{equation}
\Delta P_{\mathrm{splitter}}
= K_{\mathrm{splitter}} \frac{\rho u_c^2}{2}.
\end{equation}

Accordingly, the modeled distribution-network contribution becomes
\begin{equation}
\Delta P_{\mathrm{dist}}
=
\Delta P_{\mathrm{header}}
+ \Delta P_{\mathrm{header\rightarrow splitter}}
+ \Delta P_{\mathrm{splitter}}
+ \Delta P_{\mathrm{connector}}
+ \Delta P_{\mathrm{connector\rightarrow branch}}.
\end{equation}
The total pressure drop reported for a candidate point is the sum of these
contributions.

\subsubsection{Two-Node Bed Solve and Duty Calculation}

For each candidate operating point, the total heat delivered to the bed is
split between the bottom and top nodes using the bottom-fraction parameter
$f_b$:
\begin{equation}
Q_b = f_b Q_{\mathrm{bed,total}},
\qquad
Q_t = (1-f_b)Q_{\mathrm{bed,total}}.
\end{equation}
The steady two-node bed temperatures are then obtained from
\begin{equation}
\begin{bmatrix}
UA_{b,\infty} + UA_{bt} & -UA_{bt} \\
-UA_{bt} & UA_{t,\infty} + UA_{bt}
\end{bmatrix}
\begin{bmatrix}
T_b \\ T_t
\end{bmatrix}
=
\begin{bmatrix}
Q_b + UA_{b,\infty} T_\infty \\
Q_t + UA_{t,\infty} T_\infty
\end{bmatrix}.
\end{equation}
This $2\times2$ linear system is solved directly at every candidate point.
The required full-power duty fraction is then
\begin{equation}
\mathrm{dutyNeeded}
= \frac{Q_{\mathrm{req}}}{Q_{\mathrm{bed,total}}}.
\end{equation}

\subsubsection{Sweep Structure and Warm Starts}

For the main design sweep, the code loops over pitch, total flow, and
voltage. Within a fixed pitch and flow row, the wire-temperature iteration is
warm-started from the previous voltage point. This reduces the iteration
count because neighboring voltage states tend to have similar mean wire
temperatures.

The wire-diameter study uses a targeted subset rather than the entire design
grid. For each sampled wire diameter, the code:
\begin{enumerate}
\item fixes the comparison-selected pitch;
\item estimates a diameter-scaled voltage center from the comparison-point
      branch current and the trial-diameter branch resistance at the reference
      material state, then searches a local voltage window around that center;
\item uses the configured full-range flow subgrid for the wire study rather
      than a narrow local flow window; and
\item ranks the resulting candidate set with the same lexicographic rule used
      everywhere else in the study.
\end{enumerate}
This adaptive voltage-centering step is important for physical consistency:
when wire diameter changes, branch resistance changes strongly, so a fixed
voltage window around the original reference point can miss the relevant
operating neighborhood for thick or thin wires.

\subsubsection{Interpretation of the Moisture Maps}

The moisture figure is therefore a map of model outputs, not independent
inputs. At any point on the $(V,\dot{V}_{\mathrm{tot}})$ plane:
\begin{enumerate}
\item the imposed total inlet flow determines the inlet mass flow;
\item the heater solution determines the local heated-air temperature;
\item the heated-air density determines the released volumetric flow and
      perforation exit velocity;
\item the exit velocity and air properties determine the hole-scale
      convection and mass-transfer coefficients; and
\item the minimum of the diffusion, carrying-capacity, and heat-limited
      evaporation rates determines the water-loss and latent-load values.
\end{enumerate}
This is why all three moisture-related subplots vary with voltage, even
though the horizontal axis is not itself a flow variable.

\subsection{Standalone Summer Spot-Cooler / Assist-Blower Mode}

The JSON file now also supports a second operating mode,
\texttt{summer\_spot\_cooler}. In this mode the Python workflow does
\emph{not} call the MATLAB heater export. Instead, it constructs a
steady no-heater airflow sweep directly from the same shared wall,
aeration-tube, evaporation, and two-node bed equations already documented
above. The intent is to estimate how much airflow, spot-cooler duty, assist-blower duty, and
substrate water loss are required when the outside air is hot and the
aeration system is used as a plenum-fed cooled-air network rather than as an
electrical heater. In the current implementation, the Whynter ARC-1030WN
sets the plenum supply temperature and the downstream Goorui
GHBH 1D7 34 1R5 side-channel blower provides the pressure rise needed to move
air through the splitter and aeration tubes.

\subsubsection{Summer Boundary Conditions}

In the summer mode, ambient air first passes through a spot cooler and enters
an insulated plenum. A separate assist blower then draws from that plenum and
pushes the cooled air through the aeration network. If the
configured target supply temperature is $T_{\mathrm{supply,target}}$, then the
requested sensible cooling duty is
\begin{equation}
Q_{\mathrm{spot,req}}
= \dot m_{\mathrm{air}} c_{p,\mathrm{air}}
\max\!\left(T_{\infty,\mathrm{summer}} - T_{\mathrm{supply,target}},0\right).
\end{equation}
If the spot cooler has a configured maximum sensible capacity
$Q_{\mathrm{spot,max}}$, the applied unit duty is
\begin{equation}
Q_{\mathrm{spot}}
= \min\!\left(Q_{\mathrm{spot,req}},Q_{\mathrm{spot,max}}\right),
\end{equation}
and the actual supply temperature becomes
\begin{equation}
T_{\mathrm{air,in}}
= T_{\infty,\mathrm{summer}}
- \frac{Q_{\mathrm{spot}}}{\dot m_{\mathrm{air}} c_{p,\mathrm{air}}}.
\end{equation}
If a coefficient of performance $\mathrm{COP}$ is provided, the corresponding
spot-cooler electrical power is
\begin{equation}
P_{\mathrm{spot}} = \frac{Q_{\mathrm{spot}}}{\mathrm{COP}}.
\end{equation}
Thus the inlet air to the aeration network is no longer taken directly from
the ambient condition; it is the spot-cooled supply state:
\begin{equation}
T_{\mathrm{air,in}} = T_{\mathrm{supply,actual}},
\end{equation}
and the electrical heater is disabled, so
\begin{equation}
P_{\mathrm{tot}} = 0,
\qquad
I_{\mathrm{tot}} = 0.
\end{equation}
The assist blower is modeled by a linear fan-curve approximation. If the
published rated-duty point is
\begin{equation}
\left(\dot V_{\mathrm{rated}}, \Delta P_{\mathrm{rated}}\right),
\end{equation}
and the published shutoff pressure is
\begin{equation}
\Delta P_{\mathrm{shutoff}},
\end{equation}
then the implied free-air flow is
\begin{equation}
\dot V_{\mathrm{free}}
=
\dot V_{\mathrm{rated}}
\frac{\Delta P_{\mathrm{shutoff}}}
{\Delta P_{\mathrm{shutoff}}-\Delta P_{\mathrm{rated}}}.
\end{equation}
The available blower static at a requested airflow $\dot V$ is then
\begin{equation}
\Delta P_{\mathrm{avail}}(\dot V)
=
\max\!\left[
\Delta P_{\mathrm{shutoff}}
\left(1-\frac{\dot V}{\dot V_{\mathrm{free}}}\right),0
\right].
\end{equation}
For the current Goorui GHBH 1D7 34 1R5 basis, the model uses the listing headline values as
a temporary two-point fan-curve approximation,
\begin{equation}
\dot V_{\mathrm{rated}} = 255 \ \mathrm{m^3/h} \approx 150.1 \ \mathrm{CFM},
\qquad
\Delta P_{\mathrm{rated}} = 0 \ \mathrm{in.\ H_2O},
\qquad
\Delta P_{\mathrm{shutoff}} = 140 \ \mathrm{mbar} \approx 56.2 \ \mathrm{in.\ H_2O}.
\end{equation}
The flow sweep is therefore one-dimensional:
\begin{equation}
\dot V_{\mathrm{tot}} \in
\left[\dot V_{\min}, \dot V_{\max}\right],
\end{equation}
with the number of sampled flow points set by the JSON field
\texttt{cooling\_mode.totalFlow\_Lpm\_count}.

\subsubsection{Aeration-Tube Marching Without Electrical Heating}

The summer aeration-tube march uses the same segment structure as the heated
model, but the heater section is omitted. For each aeration segment, the
wall-transfer and jet terms remain
\begin{equation}
q_{\mathrm{wall,seg}}
= U_{\mathrm{aer}} P_{\mathrm{out}} \Delta x
\left(T_{\mathrm{air}} - T_{\mathrm{bed,ref}}\right),
\end{equation}
\begin{equation}
q_{\mathrm{jet,seg}}
= \dot m_{\mathrm{rel,seg}} c_{p,\mathrm{air}}
\left(T_{\mathrm{air}} - T_{\mathrm{bed,ref}}\right),
\end{equation}
while the latent sink is still
\begin{equation}
q_{\mathrm{evap,seg}}
= \dot m_{\mathrm{evap,seg}} h_{fg}.
\end{equation}
Accordingly, the net segment contribution to the bed is
\begin{equation}
q_{\mathrm{bed,seg}}
= q_{\mathrm{wall,seg}} + q_{\mathrm{jet,seg}} - q_{\mathrm{evap,seg}}.
\end{equation}
Summing along the tube gives
\begin{equation}
Q_{\mathrm{bed,branch}}
= \sum_{\mathrm{seg}} q_{\mathrm{bed,seg}},
\qquad
Q_{\mathrm{bed,total}} = N_{\mathrm{tubes}} Q_{\mathrm{bed,branch}}.
\end{equation}
In summer mode, net cooling corresponds to
\begin{equation}
Q_{\mathrm{bed,total}} < 0.
\end{equation}
The reported cooling capacity is therefore
\begin{equation}
Q_{\mathrm{cool}} = \max\!\left(-Q_{\mathrm{bed,total}},0\right).
\end{equation}

\subsubsection{Required Cooling Load}

The same lumped ambient-loss model is reused, but with the summer ambient
temperature replacing the winter greenhouse value. If the desired bed
reference temperature is $T_{\mathrm{set,cool}}$, the corresponding lumped
cooling requirement is
\begin{equation}
Q_{\mathrm{req,cool}}
= UA\left(T_{\infty,\mathrm{summer}} - T_{\mathrm{set,cool}}\right).
\end{equation}
This is the cooling analog of the winter heating-loss requirement and is
reported in the standalone summary as a reference capacity target.

\subsubsection{Summer Water Requirement}

The same segmentwise evaporation model is reused in summer mode, with the
inlet relative humidity taken from
\texttt{cooling\_mode.inletRelativeHumidity}. The total evaporated water is
\begin{equation}
\dot m_{\mathrm{water,day}}
= 86400 \,
N_{\mathrm{tubes}}
\sum_{\mathrm{seg}} \dot m_{\mathrm{evap,seg}}.
\end{equation}
In the current plenum-fed spot-cooler implementation, the reported quantity is the
predicted substrate water loss itself rather than a separate external
sprinkler-supply multiplier.

\subsubsection{Summer Lexicographic Selection Rule}

The summer mode also avoids weighted penalties. Instead, each point is
ranked lexicographically using the configured order in
\texttt{cooling\_mode.optimization.priority\_order}. The default hard-safe
residuals are:
\begin{align}
C_1 &= \max(\Delta P - \Delta P_{\mathrm{avail}}(\dot V),0), \\
C_2 &= \max(\Delta P - \Delta P_{\max},0), \\
C_3 &= \max(\dot m_{\mathrm{water,day}} - \dot m_{\mathrm{water,max}},0), \\
C_4 &= \max(u_h - u_{h,\max},0), \\
C_5 &= \max(T_b - T_{b,\max},0), \\
C_6 &= \max(T_t - T_{t,\max},0), \\
C_7 &= \max(|T_b-T_t| - \Delta T_{\max},0).
\end{align}
Among points for which all seven residuals are zero and
$Q_{\mathrm{bed,total}}<0$, the default clean ordering is:
\begin{enumerate}
\item minimize $\max(T_b,T_t)$;
\item minimize the mean bed temperature $\tfrac{1}{2}(T_b+T_t)$;
\item minimize combined cooling-system electrical power,
      $P_{\mathrm{spot}} + P_{\mathrm{blower}}$;
\item minimize total airflow; and
\item minimize bottom--top spread.
\end{enumerate}
Thus the summer recommendation is based directly on thermal safety,
spot-cooler duty, substrate water loss, and blower demand rather than on a
weighted penalty score.

\subsubsection{Mode Switching and Export Behavior}

The JSON runner now uses \texttt{run\_modes} to control whether the heating
branch, the summer-cooling branch, or both are executed in a single call:
\begin{enumerate}
\item \texttt{run\_modes.heating = true}: use the MATLAB-backed
      voltage-flow heater sweep and write its outputs to
      \texttt{outputs/heating};
\item \texttt{run\_modes.cooling = true}: skip the MATLAB heater export and
      build the airflow-only summer sweep directly in Python, writing its
      outputs to \texttt{outputs/cooling}; and
\item the root file \texttt{outputs/study\_summary.txt} serves as an index
      pointing to the mode-specific summaries.
\end{enumerate}
The legacy \texttt{operating\_mode} field is retained only as a
backward-compatible fallback if \texttt{run\_modes} is not present.
This separation is deliberate. The summer mode is physically the same bed,
wall, perforation, moisture, and plume model, but it is not a voltage-driven
electrical-heater problem, so it should not be represented by forcing the
heater sweep to operate at $V=0$.


\subsection{Configuration and JSON-Controlled Inputs}

The current workflow uses a shared configuration file,
\texttt{study\_config.json}, to control both the MATLAB export inputs and the
standalone Python reporting logic.

\subsubsection{Model Overrides That Change the Physics}

Edits under \texttt{model\_overrides} are passed into the MATLAB design model
before the export is generated. These settings alter the computed operating
points and therefore change the underlying physics results. The most important
groups are:
\begin{itemize}
\item \texttt{model\_overrides.sweep}, which defines the voltage, total-flow,
      pitch, and wire-diameter sweeps, including either explicit arrays or
      \texttt{min/max/count} specifications, as well as the editable report
      duty target \texttt{targetDuty}, now set to $0.75$ by default;
\item \texttt{model\_overrides.limits}, which defines explicit limits such as
      maximum current, wire temperature, outlet-air temperature, pressure
      drop, and water loss;
\item \texttt{model\_overrides.optimization}, which defines the MATLAB-side
      lexicographic thermal thresholds such as maximum bottom--top spread and
      minimum bottom and top temperatures;
\item \texttt{model\_overrides.environment}, which exposes the shared heating
      ambient temperature, ambient pressure, and reference plate length used
      in the external natural-convection estimate;
\item \texttt{model\_overrides.bin}, which now selects the active vessel-top
      condition for the export, including the closed-top vented case through
      \texttt{openTop = false}, \texttt{ventHoleCount}, and
      \texttt{ventHoleDiameter\_m}, while also exposing fill fraction, bulk
      density, bed heat capacity, effective bed conductivity, and the
      bottom-node heat-input split;
\item \texttt{model\_overrides.bin\_wall}, which contains the editable vessel
      wall-stack properties such as sheet thickness, sheet conductivity,
      insulation thickness, insulation conductivity, and the internal
      bed-to-wall convection coefficient;
\item \texttt{model\_overrides.heater\_tube}, which contains the editable
      base heater-tube dimensions such as tube length, ID, OD, and heater
      insulation thickness, as well as the heater-wall conductivity,
      heater-insulation conductivity, external convection coefficient,
      segment count, and extra heater minor-loss coefficient;
\item \texttt{model\_overrides.heaterTube}, which contains the editable
      helical-geometry inputs such as \texttt{coilMeanD\_m} and
      \texttt{coilSpan\_m};
\item \texttt{model\_overrides.aeration}, which includes the effective
      bed-side sensible coefficient $h_{\text{bed}}$, the perforation
      release fraction, the active tube count \texttt{nParallelTubes}, and
      the editable aeration-network dimensions such as length, ID, OD, and
      optional header ID, together with wall conductivity, segment count,
      the end/header minor-loss coefficients, and the explicit
      splitter/connector inputs
      \texttt{headerEnabled}, \texttt{splitterOutletCount},
      \texttt{splitterInletID\_mm}, \texttt{branchConnectorID\_mm},
      \texttt{branchConnectorLength\_m}, \texttt{splitterBodyK}, and
      \texttt{contractionConeAngle\_deg};
\item \texttt{model\_overrides.perforation}, which includes hole count and
      hole diameter; and
\item \texttt{model\_overrides.evaporation}, which includes relative
      humidity, the wetted-area basis, the Table~3 particle-surface inputs,
      surface temperature, and Lewis factor.
\end{itemize}

In the current default configuration, \texttt{headerEnabled = false}. Hence
the model does not assume a separate large upstream header; instead it
treats the distribution hardware as a bank of $4$-way splitter modules with
a nominal $19\ \mathrm{mm}$ splitter inlet and branch connector, followed by
the larger aeration tube. Contraction losses are evaluated with the
Idel'chik Section~III, Diagram~3-5 ``conical converging bellmouth without
end wall'' model at the currently configured angle of $60^\circ$. This
matches the stated use of the GroundWork
Tractor Supply $4$-way brass shutoff manifold as the physical distribution
device. Because the vendor page does not publish an internal hydraulic
diameter or an internal branch-loss coefficient, the unresolved internal
manifold loss is represented by the editable lumped parameter
\texttt{splitterBodyK}.

The \texttt{summary\_inputs} block is intentionally \emph{not} the editable
source for these geometry/material quantities. It is used only for static
report/specification text when no corresponding physics override is supplied.
Accordingly, vessel-wall properties now live under
\texttt{model\_overrides.bin\_wall}, base heater-tube dimensions under
\texttt{model\_overrides.heater\_tube}, aeration-network dimensions under
\texttt{model\_overrides.aeration}, and coil mean diameter / coil span under
\texttt{model\_overrides.heaterTube}.

If any of these quantities are changed, the export must be regenerated
because the candidate operating points themselves will change. The current
Python workflow can detect this automatically by hashing the
\texttt{model\_overrides} block and rerunning the MATLAB export when that
signature changes.

\subsubsection{Standalone Ranking and Reporting Controls}

At the top level of the same JSON file, the \texttt{optimization} block
controls the standalone Python lexicographic ranking and report display, and
the \texttt{radial\_profile} block controls the standalone worm-safe-distance
plot. The \texttt{tube\_layout} block controls the standalone cross-sectional
tube-placement, spacing, and overlap-aware worm-habitat exclusion-volume
calculation. The same file also now contains \texttt{run\_modes}, which
independently enables the MATLAB-backed heater-design branch and the
standalone summer airflow-only branch, together with the corresponding
\texttt{cooling\_mode} block. The older \texttt{operating\_mode} field is
retained only as a backward-compatible fallback if \texttt{run\_modes} is
omitted. In particular:
\begin{itemize}
\item \texttt{priority\_order} sets the ordering of the lexicographic
      criteria used in the standalone report;
\item \texttt{spread\_limit\_C}, \texttt{min\_bottom\_temp\_C}, and
      \texttt{min\_top\_temp\_C} set the corresponding displayed thermal
      thresholds; and
\item \texttt{top\_candidate\_count} and \texttt{n\_top} control how many
      ranked points are printed; and
\item \texttt{cooling\_mode.ambientAir\_C},
      \texttt{cooling\_mode.pressure\_Pa},
      \texttt{cooling\_mode.totalFlow\_Lpm\_*},
      \texttt{cooling\_mode.inletRelativeHumidity},
      \texttt{cooling\_mode.surfaceTemp\_C},
      \texttt{cooling\_mode.spot\_cooler},
      \texttt{cooling\_mode.assist\_blower}, and
      \texttt{cooling\_mode.optimization} define the standalone summer
      plenum-fed spot-cooler / assist-blower sweep, the summer hard constraints, and the
      cooling-mode lexicographic objective; and
\item \texttt{run\_modes.heating} and \texttt{run\_modes.cooling} determine
      whether the script writes the corresponding mode-specific folders
      \texttt{outputs/heating} and \texttt{outputs/cooling}, while the root
      file \texttt{outputs/study\_summary.txt} acts as an index pointing to
      the two mode-specific summaries; and
\item \texttt{radial\_profile.temperature\_threshold\_C},
      \texttt{max\_distance\_m}, \texttt{bed\_reference},
      \texttt{initial\_plume\_radius\_mode},
      \texttt{spread\_rate\_m\_per\_m}, and
      \texttt{max\_plume\_radius\_m} control the post-processed finite-plume
      radial air-cooling calculation with local re-solved evaporation; and
\item \texttt{tube\_layout.layoutMode},
      \texttt{tube\_layout.maxSupportedTubes},
      \texttt{tube\_layout.wallOffsetWidthFraction},
      \texttt{tube\_layout.wallOffsetHeightFraction},
      \texttt{tube\_layout.includeWallConduction},
      \texttt{tube\_layout.tubeCenters\_m},
      \texttt{tube\_layout.tubeCenterHeight\_m},
      \texttt{tube\_layout.tubeCenterX\_m}, and
      \texttt{tube\_layout.cross\_section\_grid\_points} control the
      cross-sectional tube spacing, the requested rule-based 1-to-12 tube
      orientation, the optional wall-return correction, and the numerical
      union-area estimate used to report pipe-occupied, worm-unsafe, and
      remaining habitable bed volume.
\end{itemize}

Within \texttt{tube\_layout}, the precedence is now explicit and consistent:
\texttt{tubeCenters\_m} fully overrides the other placement fields; if that
list is empty and \texttt{tubeCenterX\_m} is provided, the code uses a
single-row layout at \texttt{tubeCenterHeight\_m}; and if
\texttt{tubeCenterHeight\_m} is left null, the code defaults it to the same
\texttt{wallOffsetHeightFraction} times filled-bed height used by the
rule-based bottom row.

Changing only these top-level ranking or wording settings does not require a
new MATLAB export if the underlying data file is unchanged.

The current clean ordering is:
\begin{enumerate}
\item hard constraints first (current, wire temperature, outlet-air
      temperature, pressure drop, water loss, minimum bottom temperature,
      minimum top temperature, and maximum spread);
\item maximize $\min(T_b,T_t)$ among the hard-safe set;
\item minimize total electrical power;
\item minimize $|T_b-T_t|$; and
\item use mean bed temperature and residual heat shortfall only as late
      tie-breakers.
\end{enumerate}

\subsubsection{Practical Interpretation}

The JSON structure therefore separates two kinds of edits:
\begin{enumerate}
\item \textbf{Physics edits:} changes to \texttt{model\_overrides} modify the
      actual solved model and require a fresh export.
\item \textbf{Post-processing edits:} changes to top-level plot, wording, and
      ranking settings modify how the current exported dataset is interpreted
      and displayed.
\item \textbf{Summer-mode edits:} changes inside \texttt{cooling\_mode} alter
      the standalone summer plenum-fed spot-cooler / assist-blower sweep directly and do not
      require a MATLAB rerun because no heater-export dataset is used in that
      mode.
\item \textbf{Run-selection edits:} changes to \texttt{run\_modes} alter only
      which branches are executed and which output folders are written; they
      do not alter the underlying physics of either branch.
\end{enumerate}

This distinction is especially important for evaporation. Because the model
now uses
\begin{equation}
Q_{\text{bed,branch}}
=
\sum\left(q_{\text{wall}}+q_{\text{jet}}-q_{\text{evap}}\right),
\end{equation}
changing evaporation-related JSON inputs is a \emph{physics} change, not a
pure reporting change.

\subsection{Implementation Reference Map}

The MATLAB source comments use the following shorthand labels for the main
correlation families and source materials cited throughout this document:
\begin{itemize}
\item \textbf{[R1]} Churchill, S.W.; Bernstein, M. ``A Correlating Equation
      for Forced Convection from Gases and Liquids to a Circular Cylinder in
      Crossflow.'' \emph{Journal of Heat Transfer} 99(2), 1977.
\item \textbf{[R2]} Gnielinski, V. ``New Equations for Heat and Mass Transfer
      in Turbulent Pipe and Channel Flow.'' \emph{International Chemical
      Engineering} 16(2), 1976.
\item \textbf{[R3]} Churchill, S.W. ``Friction-Factor Equation Spans All
      Fluid-Flow Regimes.'' \emph{Chemical Engineering} 84(24), 1977.
\item \textbf{[R4]} Churchill, S.W.; Chu, H.H.S. ``Correlating Equations for
      Laminar and Turbulent Free Convection from a Vertical Plate.''
      \emph{International Journal of Heat and Mass Transfer} 18(11), 1975.
\item \textbf{[R5]} Bergman, Lavine, Incropera, DeWitt.
      \emph{Fundamentals of Heat and Mass Transfer}. Used for thermal
      resistance networks, continuity, energy balances, and lumped-capacitance
      arguments.
\item \textbf{[R6]} Idel'chik, I.E. \emph{Handbook of Hydraulic Resistance},
      Section~IV sudden-expansion loss, written on the upstream-velocity
      basis. This is the same Borda--Carnot coefficient presented in
      standard fluid-mechanics texts such as White or Munson.
\item \textbf{[R7]} Ideal-gas density and Sutherland-law viscosity relations
      for air.
\item \textbf{[R8]} Hartenstein, R. ``Metabolic parameters of the earthworm
      \emph{Eisenia foetida} in relation to temperature.'' \emph{Biotechnology
      and Bioengineering} 24(8), 1982.
\item \textbf{[R9]} IAPWS SR1-86(1992). Revised Supplementary Release on
      Saturation Properties of Ordinary Water Substance.
\item \textbf{[R10]} Hasan, A.; Siren, K. ``The Lewis factor and its
      influence on the performance prediction of wet-cooling towers.''
      \emph{International Journal of Thermal Sciences} 44(9), 2005.
\item \textbf{[R11]} Frederickson, J.; Howell, G.; Hobson, A.M.
      ``Effect of pre-composting and vermicomposting on compost
      characteristics.'' \emph{European Journal of Soil Biology} 43 (2007)
      S320--S326. Table~3 supplies the particle-size distribution used in
      the current wetted-area estimate.
\end{itemize}

\documentclass{article}
\usepackage{amsmath}
\usepackage[margin=1in]{geometry}

\subsection{Year-Round Energy, Cost, and Climate-Driven Operation}

This note documents the Python year-round layer that extends the main
vermicomposter heating/cooling study. The annual layer is not an independent
thermal model. It reuses:
\begin{enumerate}
\item the current two-node bed model,
\item the currently selected heating reference point from the MATLAB-backed
heating study, and
\item the currently selected cooling reference point from the standalone
summer spot-cooler plus assist-blower study.
\end{enumerate}

The new annual layer differs from the earlier monthly
\texttt{resistiveheatingcode.m} scaffold in one important way: the ambient
driver is now a source-backed representative Connecticut climate year rather
than a short hard-coded monthly average list plus ad hoc offsets.

\section{Climate Data Sources}

\subsection{Connecticut Monthly Mean and Variability}

Monthly ambient statistics are taken from NOAA NCEI Climate at a Glance
(CAG), using the statewide Connecticut monthly mean temperature series
(\texttt{tavg}) for calendar years 2016--2025. The CAG API is documented at
\begin{equation}
\texttt{https://www.ncei.noaa.gov/monitoring-content/cag/text/time-series-api-doc.html}
\end{equation}
and the monthly statewide endpoint pattern used in the study is
\begin{equation}
\texttt{https://www.ncei.noaa.gov/access/monitoring/climate-at-a-glance/}
\end{equation}
\begin{equation}
\texttt{statewide/time-series/6/tavg/1/\{month\}/2016-2025/data.json}
\end{equation}
where state code $6$ is Connecticut.

From that NOAA series, the Python model stores for each month $m$:
\begin{equation}
\bar{T}_m
\end{equation}
as the 2016--2025 monthly mean temperature in $^\circ$C, and
\begin{equation}
\sigma_m
\end{equation}
as the corresponding 2016--2025 monthly standard deviation of that monthly
mean series.

\subsection{Freeze and Heat-Wave Timing Statistics}

Event timing is derived from NOAA daily-summaries station
\texttt{USW00014740} over the same 2016--2025 period using
\begin{equation}
\texttt{https://www.ncei.noaa.gov/access/services/data/v1?dataset=daily-summaries}
\end{equation}
with daily \texttt{TMAX}, \texttt{TMIN}, and \texttt{TAVG}.

The annual model uses this daily dataset only to derive event timing and
frequency statistics. The ambient baseline itself remains the Connecticut
statewide monthly mean series.

\section{Representative-Year Construction}

\subsection{Gaussian-Regression Baseline Profile}

The smooth daily baseline is now built by periodic Gaussian kernel
regression, not by a sinusoid. Let the month-midpoint day-of-year anchors be
\begin{equation}
x_m,
\qquad m=1,\dots,12,
\end{equation}
with observed monthly mean temperatures
\begin{equation}
\bar{T}_m.
\end{equation}
For any day of year $d$, the circular day-distance to month anchor $m$ is
\begin{equation}
\delta(d,x_m)
=
\min\!\left(|d-x_m|,\ 365-|d-x_m|\right).
\end{equation}
With Gaussian kernel bandwidth $h$, the daily baseline ambient is
\begin{equation}
T_{\mathrm{base}}(d)
=
\frac{
\sum_{m=1}^{12}
\exp\!\left(-\delta(d,x_m)^2/(2h^2)\right)\bar{T}_m
}{
\sum_{m=1}^{12}
\exp\!\left(-\delta(d,x_m)^2/(2h^2)\right)
}.
\end{equation}

The same periodic Gaussian regression is also applied to the monthly standard
deviation anchors $\sigma_m$ to build a smooth daily variability envelope,
although the primary annual ambient baseline is the regressed mean field
$T_{\mathrm{base}}(d)$.

This change was made because the earlier sine-based interpolation imposed an
artificial oscillation within every month. The Gaussian-kernel reconstruction
is a smoother nonparametric interpolation of the monthly Connecticut climate
anchors and is therefore a more defensible baseline for the representative
year.

\subsection{Freeze-Day Definition and Override}

The study uses the operational freeze-day definition
\begin{equation}
T_{\min} \le 32^\circ\mathrm{F},
\end{equation}
applied to the NOAA daily-summaries record.

For each month, the daily-station analysis provides:
\begin{enumerate}
\item the average number of freeze days per month,
\item the average day of month on which freeze days occur, and
\item the mean freeze-day anomaly below the corresponding Connecticut monthly
mean baseline.
\end{enumerate}

Let the expected freeze-day count for month $m$ be $F_m$. The code converts
the expected monthly counts to a representative integer annual schedule using
largest-remainder allocation:
\begin{equation}
n_{f,m} = \mathrm{LR}(F_m),
\end{equation}
where the total number of allocated freeze days equals the rounded annual
expected total.

The freeze days are \emph{not} then forced into one contiguous block. Instead,
the code assigns each day-of-month $i$ a Gaussian timing weight centered on
the mean observed freeze day $\mu_{f,m}$:
\begin{equation}
w_{f,m}(i)
=
\exp\!\left(
-\frac{(i-\mu_{f,m})^2}{2s_{f,m}^2}
\right),
\end{equation}
where $s_{f,m}$ is a month-scale spread parameter. Freeze days are then
selected sequentially by maximizing a repulsion-modified score
\begin{equation}
\tilde{w}_{f,m}(i)
=
w_{f,m}(i)
\prod_{j\in\mathcal{S}_{f,m}}
\left[
1-0.65\exp\!\left(
-\frac{(i-j)^2}{2r_{f,m}^2}
\right)
\right],
\end{equation}
with $\mathcal{S}_{f,m}$ the set of already selected freeze days and
$r_{f,m}$ a repulsion length scale. This preserves the NOAA-derived mean timing
but discourages selecting many neighboring days, so the representative freeze
days are dispersed through the cold part of the month rather than forming one
artificial mid-month block.

If the mean freeze anomaly is $\Delta T_{f,m}>0$, the ambient profile is
overridden by
\begin{equation}
T_m(d)
=
\min\!\left(
T_{\mathrm{base},m}(d),
\bar{T}_m - \Delta T_{f,m}
\right)
\end{equation}
on the selected freeze days.

\subsection{Heat-Wave Definition and Override}

The study uses the operational heat-wave definition
\begin{equation}
\text{at least 3 consecutive days with } T_{\max} \ge 90^\circ\mathrm{F}.
\end{equation}

This threshold is a study definition applied consistently to the NOAA daily
record. From the daily station data, the code derives for each month:
\begin{enumerate}
\item average heat-wave spell starts per month,
\item average heat-wave days per month,
\item average heat-wave start day of month, and
\item mean heat-wave anomaly above the Connecticut monthly baseline.
\end{enumerate}

If $S_m$ is the expected monthly spell-start count and $H_m$ is the expected
monthly heat-wave day count, the code allocates representative integer counts
by largest-remainder rounding:
\begin{equation}
n_{s,m} = \mathrm{LR}(S_m),
\qquad
n_{h,m} = \mathrm{LR}(H_m).
\end{equation}

Within each month, the spell model enforces the study definition by requiring
at least three days per spell. The spell-start candidates are selected from a
Gaussian timing distribution centered on the observed mean heat-wave start day
$\mu_{h,m}$:
\begin{equation}
w_{h,m}(i)
=
\exp\!\left(
-\frac{(i-\mu_{h,m})^2}{2s_{h,m}^2}
\right),
\end{equation}
then dispersed by the same repulsion concept used for freeze days:
\begin{equation}
\tilde{w}_{h,m}(i)
=
w_{h,m}(i)
\prod_{j\in\mathcal{S}_{h,m}}
\left[
1-0.65\exp\!\left(
-\frac{(i-j)^2}{2r_{h,m}^2}
\right)
\right].
\end{equation}
Here $\mathcal{S}_{h,m}$ contains already assigned spell starts and
$r_{h,m}$ is scaled to the representative spell spacing. After these starts are
chosen, the heat-wave days remain contiguous \emph{within each spell} because
that contiguity is part of the heat-wave definition itself. The change is only
that the spell starts are now distributed across the month instead of being
packed around a single central date.

If the mean heat-wave anomaly is $\Delta T_{h,m}>0$, the ambient profile is
overridden on those heat-wave days by
\begin{equation}
T_m(d)
=
\max\!\left(
T_m(d),
\bar{T}_m + \Delta T_{h,m}
\right).
\end{equation}

\subsection{Resulting Representative-Year Ambient Driver}

The final ambient driver is therefore a 365-day deterministic sequence
constructed from:
\begin{enumerate}
\item statewide Connecticut monthly means,
\item statewide monthly variability levels,
\item freeze-day frequency, timing, and cold anomaly from NOAA daily station
data, and
\item heat-wave spell frequency, timing, duration, and hot anomaly from NOAA
daily station data.
\end{enumerate}

\section{Daily Heating and Cooling Loads}

Let the current bin conductance to ambient be
\begin{equation}
UA_{\mathrm{bin}}.
\end{equation}
For each representative day $d$ with ambient temperature
$T_\infty(d)$, the required heating load is
\begin{equation}
Q_{\mathrm{req},h}(d)
=
UA_{\mathrm{bin}}
\max\!\left(T_{h,\mathrm{set}} - T_\infty(d),\,0\right),
\end{equation}
and the required cooling load is
\begin{equation}
Q_{\mathrm{req},c}(d)
=
UA_{\mathrm{bin}}
\max\!\left(T_\infty(d) - T_{c,\mathrm{set}},\,0\right).
\end{equation}

\section{Reference-Point Capacity Application}

The annual layer no longer treats the selected heating and cooling points as
fixed-capacity, fixed-power devices evaluated only at the original design
ambients of $-15^\circ\mathrm{C}$ and $37.5^\circ\mathrm{C}$. Instead, it
keeps the selected operating \emph{controls} fixed and re-solves those same
operating points at each representative day's ambient temperature.

For heating, the fixed controls are the selected branch voltage, total flow,
pitch, wire geometry, and current tube/manifold geometry. At each day $d$
with ambient temperature $T_\infty(d)$, the Python annual layer re-solves the
same heater-tube and aeration-tube balances used by the design model to
obtain an ambient-specific heating capacity and electrical draw:
\begin{equation}
Q_{h,\mathrm{cap}}(d)
=
Q_{\mathrm{to\,bed},h}\!\left(T_\infty(d)\right),
\end{equation}
\begin{equation}
P_{h,\mathrm{elec}}(d)
=
P_h\!\left(T_\infty(d)\right).
\end{equation}

For cooling, the fixed controls are the selected total flow, spot-cooler
settings, and assist-blower configuration. The annual layer then re-solves
the standalone cooling point at the day's ambient temperature to obtain
\begin{equation}
Q_{c,\mathrm{cap}}(d)
=
-Q_{\mathrm{to\,bed},c}\!\left(T_\infty(d)\right),
\end{equation}
with total electrical input
\begin{equation}
P_{c,\mathrm{elec}}(d)
=
P_{\mathrm{spot\,cooler}}(d) + P_{\mathrm{assist\,blower}}(d).
\end{equation}

For each day, the required duty fractions are
\begin{equation}
D_h^\ast(d) = \frac{Q_{\mathrm{req},h}(d)}{Q_{h,\mathrm{cap}}(d)},
\qquad
D_c^\ast(d) = \frac{Q_{\mathrm{req},c}(d)}{Q_{c,\mathrm{cap}}(d)},
\end{equation}
and the applied duties are clipped at unity:
\begin{equation}
D_h(d)=\min(\max(D_h^\ast(d),0),1),
\qquad
D_c(d)=\min(\max(D_c^\ast(d),0),1).
\end{equation}

If the day is a heating day, the applied bed load is
\begin{equation}
Q_{\mathrm{bed}}(d) = D_h(d)\,Q_{h,\mathrm{cap}}(d),
\end{equation}
whereas on a cooling day it is
\begin{equation}
Q_{\mathrm{bed}}(d) = -D_c(d)\,Q_{c,\mathrm{cap}}(d).
\end{equation}

\section{Daily Bottom and Top Bed Temperatures}

The current two-node steady model is then solved each day using the same bed
split fraction $f_b$ as the main study:
\begin{equation}
\begin{bmatrix}
UA_{b,\infty} + UA_{bt} & -UA_{bt} \\
-UA_{bt} & UA_{t,\infty} + UA_{bt}
\end{bmatrix}
\begin{bmatrix}
T_b(d) \\ T_t(d)
\end{bmatrix}
=
\begin{bmatrix}
f_b Q_{\mathrm{bed}}(d) + UA_{b,\infty}T_\infty(d) \\
(1-f_b) Q_{\mathrm{bed}}(d) + UA_{t,\infty}T_\infty(d)
\end{bmatrix}.
\end{equation}

Therefore the annual temperature plot is not heuristic. It is the same bed
energy-balance model already used in the heating and cooling studies, solved
day by day under the representative Connecticut climate driver.

\section{Daily Electricity, Cost, and Water}

Daily electricity use is
\begin{equation}
E_h(d) = \frac{P_{h,\mathrm{elec}}(d)\, D_h(d) \cdot 24}{1000},
\qquad
E_c(d) = \frac{P_{c,\mathrm{elec}}(d)\, D_c(d) \cdot 24}{1000},
\end{equation}
in $\mathrm{kWh/day}$. Total daily electricity use is
\begin{equation}
E(d) = E_h(d) + E_c(d).
\end{equation}

This matters because the annual electricity model now reflects the fact that
heater-shell losses, air density, wire resistance, spot-cooler sensible load,
and assist-blower duty all vary with the day's ambient condition. The annual
cost is therefore no longer based on simply scaling the single design-point
power values found at $-15^\circ\mathrm{C}$ and $37.5^\circ\mathrm{C}$.

If the Connecticut electricity price is $c_{\mathrm{elec}}$ in
USD/kWh, then daily operating cost is
\begin{equation}
C(d) = c_{\mathrm{elec}} E(d).
\end{equation}

The current default JSON value is
\begin{equation}
c_{\mathrm{elec}} = 0.2875 \ \mathrm{USD/kWh},
\end{equation}
based on the U.S. EIA Electric Power Annual 2024, Table 2.10, Connecticut
residential average retail electricity price.

The daily water replacement estimate scales the already-computed heating or
cooling water-loss rate:
\begin{equation}
m_{\mathrm{water}}(d)
=
\dot{m}_{\mathrm{water},24h}\,D(d),
\end{equation}
so the annual layer does not create a second, inconsistent moisture model.

\section{Monthly Aggregation}

After the daily solution is built, the code aggregates by month to report:
\begin{enumerate}
\item ambient minimum, maximum, mean, and source-backed monthly standard
deviation,
\item freeze-day and heat-wave counts,
\item mean and worst-case bottom/top bed temperatures,
\item mean and peak daily duty,
\item monthly heating and cooling electricity,
\item monthly operating cost, and
\item monthly water replacement.
\end{enumerate}

The time-series figure written to
\texttt{outputs/year\_round/year\_round\_energy\_cost\_analysis.png}
contains:
\begin{enumerate}
\item ambient, bottom-node, and top-node temperatures versus day of year,
\item explicit freeze-day and heat-wave day markers, and
\item daily plus cumulative operating cost versus day of year.
\end{enumerate}

\section{Relation to \texttt{resistiveheatingcode.m}}

The local MATLAB file \texttt{resistiveheatingcode.m} was used as the
starting point for the annual idea, but the current Python model is more
physically consistent because it:
\begin{enumerate}
\item removes the older ad hoc monthly solar and passive-cooling offsets,
\item replaces the short hard-coded monthly temperature list with NOAA-backed
Connecticut climate inputs, and
\item computes daily bottom and top bed temperatures from the same current
two-node model used in the rest of the study.
\end{enumerate}

\section{Sources}

\begin{enumerate}
\item NOAA NCEI Climate at a Glance API documentation:
\texttt{https://www.ncei.noaa.gov/monitoring-content/cag/text/time-series-api-doc.html}
\item NOAA NCEI Climate at a Glance Connecticut statewide monthly mean
temperature endpoint pattern:
\texttt{https://www.ncei.noaa.gov/access/monitoring/climate-at-a-glance/statewide/time-series/6/tavg/1/\{month\}/2016-2025/data.json}
\item NOAA daily-summaries endpoint used for event statistics:
\texttt{https://www.ncei.noaa.gov/access/services/data/v1?dataset=daily-summaries\&stations=USW00014740\&startDate=2016-01-01\&endDate=2025-12-31\&dataTypes=TMAX,TMIN,TAVG\&units=standard\&format=json}
\item U.S. Energy Information Administration, Electric Power Annual 2024,
Table 2.10:
\texttt{https://www.eia.gov/electricity/annual/html/epa\_02\_10.html}
\item The main study heat-transfer and moisture equations documented in
\texttt{VermiHeatLoss\_complete\_updated.tex} and
\texttt{vermi\_water\_optimization\_numerics\_appendix.tex}; the annual layer
reuses that same two-node bed model and the currently selected heating and
cooling reference points.
\end{enumerate}



\end{document}
